% -*- fill-column: 80; -*-
\documentclass[review,anonymous,screen,acmsmall,nonacm]{acmart}
\usepackage{mathpartir}
\usepackage{tikz-cd}
\usepackage{lipsum}
\usepackage{enumitem}
\usepackage{wrapfig}
\usepackage{fancyvrb}
\usepackage{xspace}

\usepackage[LGR,T1]{fontenc}
\DeclareMathAlphabet{\mathgtt}{LGR}{cmtt}{m}{n}

\renewcommand{\Gamma}{\mathgtt{G}}
\renewcommand{\Delta}{\mathgtt{D}}
\renewcommand{\Sigma}{\mathgtt{S}}
\renewcommand{\mu}{\mathgtt{m}}

\renewcommand{\familydefault}{\ttdefault}
\usepackage{mathastext}
\renewcommand{\familydefault}{\rmdefault}

\definecolor{codegreen}{rgb}{0,0.6,0}
\definecolor{codegray}{rgb}{0.5,0.5,0.5}
\definecolor{codepurple}{rgb}{0.58,0,0.82}
\definecolor{backcolour}{rgb}{0.95,0.95,0.92}

\usepackage{listings}
\lstdefinelanguage{Agda}{
keywords={data, where, module, import, open, public,
record, field, let, in, if, then, else, case, of, with,
do, postulate, primitive, mutual, abstract, private,
forall, exists, cong, set, prop, sort, Level, Data, Type,
Renamer},
morekeywords=[2]{Set, Type, Prop, moduleFuncs, instance, Named},
sensitive=true,
comment=[l]{--},
morecomment=[s]{\{-}{-\}},
morestring=[b]",
mathescape=true,
escapeinside={(*@}{@*)}
}
\lstset{
language=Agda,
basicstyle=\ttfamily\small,
keywordstyle=\color{blue},
keywordstyle=[2]\color{teal},
identifierstyle=\color{black},
commentstyle=\color{gray}\textit,
stringstyle=\color{orange},
numbers=none,
numberstyle=\tiny\color{gray},
stepnumber=1,
numbersep=5pt,
showstringspaces=false,
tabsize=4,
captionpos=b,
breaklines=true,
breakatwhitespace=false,
rulecolor=\color{black},
texcl=true,
backgroundcolor=\color{backcolour},
frame=single
}

\usepackage{hyperref}
\usepackage[capitalise]{cleveref}
\usepackage{stmaryrd}
\usepackage{scalerel}
\usepackage{tikz}
\usetikzlibrary{automata, positioning, arrows, fit}

\usepackage{pdfpages}

\newcommand{\SPF}{\texttt{SPF}}
\newcommand{\Var}{\texttt{Var}}
\newcommand{\K}{\texttt{K}}
\newcommand{\map}{\texttt{map}}
\newcommand{\roll}{\texttt{roll}}
\newcommand{\fold}{\texttt{fold}}
\newcommand{\inl}{\texttt{inl}}
\newcommand{\inr}{\texttt{inr}}
\newcommand{\sem}[1]{\llbracket{#1}\rrbracket}
\newcommand{\cat}[1]{\mathbf{#1}}
\newcommand{\lto}{\multimap}
\newcommand{\tol}{\mathrel{\rotatebox[origin=c]{180}{$\lto$}}}
\newcommand{\String}{\textbf{String}}
\newcommand{\stringg}{\texttt{String}}
\newcommand{\charg}{\mathtt{Char}}
\newcommand{\Set}{\mathbf{Set}}
\newcommand{\Syn}{\mathbf{Synx}}
\newcommand{\SemAct}{\mathbf{SemAct}}
\newcommand{\Gr}{\mathbf{Gr}}
\newcommand{\Grammar}{\mathbf{Gr}}
\newcommand{\Type}{\mathbf{Type}}
\newcommand{\Prop}{\mathbf{Prop}}
\newcommand{\Bool}{\mathtt{Bool}}
\newcommand{\true}{\mathtt{true}}
\newcommand{\false}{\mathtt{false}}
\newcommand{\nat}{\mathbb{N}}
\newcommand{\theoryname}{Dependent Lambek Calculus\xspace}
\newcommand{\theoryabbv}{$\textrm{Lambek}^D$~}
\newcommand{\lnld}{$\textrm{LNL}_D$}

\newcommand{\isTy}{\textrm{ type}}
\newcommand{\isCtx}{\textrm{ ctx}}
\newcommand{\isSmall}{\textrm{ small}}
\newcommand{\isLinTy}{\textrm{ lin. type}}
\newcommand{\isLinCtx}{\textrm{ lin. ctx.}}

\newcommand{\quoteTy}[1]{\lceil{#1}\rceil}
\newcommand{\unquoteTy}[1]{\lfloor{#1}\rfloor}

\newcommand{\gluedNL}{{\mathcal G}_S}
\newcommand{\gluedNLUniv}{{\mathcal G}_{S,i}}
\newcommand{\gluedL}{{\mathcal G}_L}

\newcommand{\amp}{\mathrel{\&}}
\newcommand{\pair}{\amp}
\DeclareMathOperator*{\bigamp}{\scalerel*{\&}{\bigoplus}}
\DeclareMathOperator*{\bigwith}{\scalerel*{\&}{\bigoplus}}


\newcommand{\bang}{~\textbf{!}~}
% Tentatively using uparrow for linear to non-linear to be in line with
% Pfennings adjoint functional programming
\newcommand{\ltonl}[1]{~\uparrow #1}
\newcommand{\nil}{\texttt{nil}}
\newcommand{\ident}{\texttt{id}}
\newcommand{\cons}{\texttt{cons}}
\newcommand{\epscons}{\varepsilon\texttt{cons}}
\newcommand{\data}{\mathsf{data}}
\newcommand{\where}{\mathsf{where}}
\newcommand{\Trace}{\texttt{Trace}}
\newcommand{\literal}[1]{\texttt{\textquotesingle#1\textquotesingle}}
\newcommand{\stringquote}[1]{\texttt{\textquotedbl#1\textquotedbl}}
\newcommand{\internalize}[1]{\lceil#1\rceil}

\newcommand{\oplusinj}[2]{\sigma\,#1\,#2}
\newcommand{\withprj}[2]{\pi\,#1\,#2}

\newcommand{\simulsubst}[2]{#1\{#2\}}
\newcommand{\subst}[3]{\simulsubst {#1} {#2/#3}}
\newcommand{\el}{\mathsf{el}}
\newcommand{\letin}[3]{\mathsf{let}\, #1 = #2 \, \mathsf{in}\, #3}
\newcommand{\lamb}[2]{\lambda #1.\, #2}
\newcommand{\lamblto}[2]{\lambda^{{\lto}} #1.\, #2}
\newcommand{\lambtol}[2]{\lambda^{{\tol}} #1.\, #2}
\newcommand{\dlamb}[2]{{\lambda}^{{\&}} #1.\, #2}
\newcommand{\withlamb}[2]{{\lambda}^{{\&}} #1.\, #2}
\newcommand{\app}[2]{#1 \, #2}
\newcommand{\applto}[2]{#1 \, #2}
\newcommand{\apptol}[2]{#1 \mathop{{}^{\tol}} #2}
\newcommand{\PiTy}[3]{\textstyle\prod (#1 : #2). #3}
\newcommand{\SigTy}[3]{\textstyle\sum (#1 : #2). #3}
\newcommand{\LinPiTy}[3]{\textstyle\bigamp (#1 : #2). #3}
\newcommand{\LinPiTyLimit}[3]{\bigwith\limits_{#1 : #2} #3}
\newcommand{\LinSigTy}[3]{\textstyle\bigoplus (#1 : #2). #3}
\newcommand{\LinSigTyLimit}[3]{\bigoplus\limits_{#1 : #2} #3}
%% \newcommand{\DepWith}[2]{{\textstyle\bigamp}\limits_{#1}{#2}
%% \newcommand{\DepPlus}[2]{{\textstyle\bigoplus}\limits_{#1}{#2}
\newcommand{\GrTy}{\mathsf{Gr}}

\newcommand{\equalizer}[3]{\{#1\,|\,\applto {#2}{#1} = \applto{#3}{#1} \}}
\newcommand{\equalizerin}[1]{\langle #1 \rangle}
\newcommand{\equalizerpi}[1]{#1.\pi}

\newcommand{\ctxwff}[1]{#1 \isCtx}
\newcommand{\ctxwffjdg}[2]{#1 \vdash #2 \isTy}
\newcommand{\linctxwff}[2]{#1 \vdash #2 \isLinCtx}
\newcommand{\linctxwffjdg}[2]{#1 \vdash #2 \isLinTy}

\newsavebox{\logoagdabox}
\sbox{\logoagdabox}{%
  %
  \raisebox{-2pt}{\includegraphics[height=1em]{logo-agda.pdf}}%
}
\newcommand{\agdalogo}{%
  \usebox{\logoagdabox}}%

\newif\ifdraft
\drafttrue
\newcommand{\steven}[1]{\ifdraft{\color{orange}[{\bf Steven says}: #1]}\fi}
\renewcommand{\max}[1]{\ifdraft{\color{blue}[{\bf Max says}: #1]}\fi}
\newcommand{\pedro}[1]{\ifdraft{\color{red}[{\bf Pedro says}: #1]}\fi}
\newcommand{\nathan}[1]{\ifdraft{\color{green}[{\bf Nathan says}: #1]}\fi}
\newcommand{\pipe}{\,|\,}

\begin{document}

\pagestyle{plain}

\pagebreak

\title{Intrinsic Verification of Parsers and Formal Grammar Theory in Dependent Lambek Calculus}

\author{Steven Schaefer}
\affiliation{\department{Electrical Engineering and Computer Science}
  \institution{University of Michigan}
  \country{USA}
}
\email{stschaef@umich.edu}

\author{Nathan Varner}
\affiliation{\department{Electrical Engineering and Computer Science}
  \institution{University of Michigan}
  \country{USA}
}
\email{nvarner@umich.edu}

\author{Max S. New}
\affiliation{
  \department{Electrical Engineering and Computer Science}
  \institution{University of Michigan}
  \country{USA}
}
\email{maxsnew@umich.edu}

\author{Pedro H. Azevedo de Amorim}
\affiliation{
  \department{Department of Computer Science}
  \institution{University of Oxford}
  \country{UK}
}
\email{pedro.azevedo.de.amorim@cs.ox.ac.uk}

\makeatletter
\let\@authorsaddresses\@empty
\makeatother

\begin{abstract}
  We propose \theoryname~(\theoryabbv), a domain-specific dependent type theory
  for the definition of intrinsically verified parsers and verification of weak
  and strong equivalences of formal grammars. The linear types in \theoryabbv
  have a simple interpretation as formal grammars, with values of the types
  interpreted as parse trees and general terms as a form of intrinsically
  verified parse transformers, of which parsers being a special case. We
  demonstrate the expressivity of this system by showing that the indexed
  inductive linear types in this calculus can be used to encode many different
  familiar grammar formalisms, including regular and context-free grammars but
  also traces of various types of automata. We can then construct strong
  equivalences between grammar and automata formalisms and construct verified
  parsers as simple linear functional programs. All of our examples have been
  implemented in a prototype implementation of \theoryabbv as a shallow
  embedding in Agda.
\end{abstract}

\maketitle

\section{Introduction}
\label{sec:intro}
\steven{This intro sounds reallchoppy but follows the following outline.

  \begin{itemize}
  \item parsing correctness is imporant because parsing is
  everywhere and failure
          could even invalidate other proof efforts.
  \item verified parsers should be
          easier to write
  \item Conal Elliott gives the right idea for what the foundation
          of parser verification should be --- type theory

  \end{itemize}
          }
Formal verification encompasses a broad suite of techniques one may use to
mathematically guarantee the correctness of software systems. We can only
reap the benefits of
any proof effort if its underlying assumptions hold. For instance, verifying the
functional correctness of an algorithm is no use if it is run on buggy hardware.
Likewise, the promised guarantees of a program may be jeopardized if it is
parsed incorrectly.

Large verified developments, like CompCert \cite{leroy_formal_2009},
invest enormous amounts of person-hours into proof development. In early
versions of CompCert, the apparent
assurances associated with their effort could indeed be invalidated through
parser bugs \cite{yangFindingUnderstandingBugs}. Parsing is a ubiquitous task
throughout computing and yet parsers are frequently unverified, even in
otherwise verified systems. The CompCert team later
developed a method for validating the correctness of LR parsers
\cite{jourdanValidatingLRParsers2012} and then replaced their parser with one
that is certified to be correct. Despite their success in  validating their C parser and other work in the
generation of  verified parsers
\cite{lasserCoStarVerifiedALL2021,firsovCertifiedNormalizationContextFree2015,danielssonTotalParserCombinators2010},
it is still arduous and highly nontrivial to write and verify a
parser. Further, distinct approaches to generating verified parsing
implementations are often relative to a different grammar formalism. We believe
that these sepearate grammar formalisms can all be unified into a
syntax-independent notion of grammar, and that this notion of grammar can be
axiomatized into a syntax that is useful for building verified parsers.

\steven{I don't like how I make the cited work sound bad above}

Elliott \cite{elliottSymbolicAutomaticDifferentiation2021}
notes that grammars can be described as type-level predicates on strings that prove
language membership. That is, a \emph{formal grammar} $A$ is a function
$\String \to \Set$ such that for a string $w$, $A~w$ is the set of proofs
showing that $w$ belongs to the language recognized by $A$. We adopt this
definition of a formal grammar.

In this paper, we present \theoryname (\theoryabbv) as a domain specific
type theory for building verified parsers. The type theory itself is a form of dependent linear-non-linear type
theory: the linear types can be interpreted as grammars, whereas the
non-linear types can be interpreted as ordinary sets. As in prior work
on dependent linear-non-linear types, the linear types are allowed to
depend on the non-linear types but not vice-versa \cite{krishnaswami_integrating_2015}.

The substructural nature of \theoryabbv is well-aligned with the
requirements intrinsic to parsing and to the theory of formal
languages, where strings constitute a very clear notion of resource
that cannot be duplicated, reordered, or dropped. Moreover, the constructive
aspect of \theoryabbv ensures that verification of parsers written in the
calculus are \emph{correct-by-construction}. The type system is powerful enough
that derivations of a term type checking carry intrinsic proofs of correctness.

\paragraph{Contributions} This paper begins in \cref{sec:type-theory-examples}
by giving small examples of programs in \theoryabbv so that we may build intuition.
From there, in \cref{sec:tt} we provide the syntax of \theoryname, a type theory
that axiomatizes the underlying structures needed to conduct formal grammar
theory. In \cref{sec:applicatoins} we demonstrate the
applicablity of \theoryabbv for defining parser terms and representing results
from formal language theory as proof terms in our calculus.
Then in \cref{sec:semantics-and-metatheory}, we give a denotational semantics
that makes precise the connection between \theoryabbv syntax and formal grammars.

Our contributions are then:
%
\begin{itemize}
  \item \theoryname (\theoryabbv): A dependent
        linear-non-linear type theory for building verified parsers
  \item An implementation of \theoryabbv, which is embedded shallowly in Cubical
        Agda. Throughout this paper, we will use \agdalogo~ to mark results that
        have been mechanized in our Agda development.
\end{itemize}

\section{\theoryname by Example}
\label{sec:type-theory-examples}
To gain intuition for working in \theoryabbv, we begin with some illustrative
examples drawn from the theory of formal languages. In particular, we
demonstrate how to encode finite grammars, regular expressions, and finite
automata in \theoryabbv. Each of our examples will be defined for strings over the three character alphabet
$\Sigma = \{ \texttt{a} , \texttt{b}, \texttt{c} \}$.

\newcommand{\A}{\texttt{A}}
\newcommand{\B}{\texttt{B}}
\newcommand{\I}{\texttt{I}}
\newcommand{\f}{\texttt{f}}
\newcommand{\g}{\texttt{g}}
\renewcommand{\L}{\texttt{L}}
\renewcommand{\a}{\texttt{a}}
\renewcommand{\b}{\texttt{b}}
\renewcommand{\c}{\texttt{c}}
\newcommand{\w}{\texttt{w}}

\paragraph{Finite Grammars}
First consider finite grammars --- those built from base types via disjunctions and
concatenations. The base types comprise characters drawn from the alphabet, the
empty string, and the empty grammar.
For instance, the grammar corresponding to the character $a$ is given by $\literal a$ which
has a single parse tree for the string
\stringquote{a} and no parse trees at any other strings. The grammar $\I$ recognizes has a single
parse tree for the empty string $\epsilon = \stringquote{}$ and no parses for any other strings.
The final base type, the empty grammar $\bot$, has no parses for any string. We
use type-theoretic syntax to represent disjunction $\oplus$ and concatenation
$\otimes$ of
grammars. Over an input string $\w$, a parse of the disjunction $\A \oplus \B$ is either
a parse of $\A$ over the string $\w$ or a
parse of $\B$ over the string $\w$. Similarly, $\w$ matches $\A \otimes \B$ if
$\w$ can be split into two strings $\w_{\A}$ and $\w_{\B}$ that match $\A$
and $\B$, respectively.

For a grammar $\A$, a parse tree of a string $\w$ can be represented as
a term of type $\A$ in the context $\internalize \w$, where
$\internalize \w$ is a context with one variable for each character of
$\w$. For example, to define a parse tree for \stringquote{ab}, we use
the context $\internalize{\stringquote{ab}} = a : \literal a , b : \literal b$. Then the left of
\cref{fig:fingram} gives a typing derivation that shows that $\inl(a
\otimes b)$ defines a valid parse tree for the finite grammar
$(\literal a \otimes \literal b)
\oplus \literal c$. Underneath the derivation in \cref{fig:fingram}, we give a linear $\lambda$-term
that implements the program given in the derivation.

\begin{figure}
\begin{mathpar}
  \inferrule
  {
    \inferrule
    {
      \inferrule
      {~}
      {a : \literal a \vdash a : \literal a}
      \\
      \inferrule
      {~}
      {b : \literal b \vdash b : \literal b}
    }
    {a : \literal a , b : \literal b \vdash a \otimes b : \literal a \otimes
      \literal b}
  }
  {a : \literal a , b : \literal b \vdash \texttt{f} := \inl(a \otimes b) :
    (\literal a \otimes \literal b) \oplus \literal c}
\end{mathpar}
\begin{lstlisting}
f : (*@$\uparrow$@*)('a' (*@$\otimes$@*) 'b' (*@$\lto$@*) ('a' (*@$\otimes$@*) 'b') (*@$\oplus$@*) 'c')
f (a , b) = inl (a (*@$\otimes$@*) b)
\end{lstlisting}
\caption{\stringquote{ac} matches
  $(\literal a \oplus \literal b) \otimes \literal c$}
\label{fig:fingram}
\end{figure}
For this interpretation of terms as parse trees to make sense, our
calculus cannot allow for \emph{any} of the usual structural rules of
a type theory: weakening, contraction and exchange. Weakening allows
for variables to go unused, while contraction allows for the same
variable to be used twice, but in a parse tree, every character must
be accounted for exactly once. That is, we want to prevent the following
erroneous derivations,
\[
  a : \literal a , b : \literal b \not \vdash a : \literal a \qquad a : \literal a , b : \literal b \not \vdash (a, a) : \literal a \otimes \literal a
\]

%
Finally, the ordering of characters in a string cannot be ignored while
parsing, so we omit the exchange rule because it would allow
for variables in the context to be reordered,
\[
  a : \literal a , b : \literal b \not \vdash (b , a) : \literal b \otimes \literal a
\]
% otherwise we would have that $\inl(x\otimes y)$ is a valid
% term in context $y:b,x:a$, but there should be no parses of $ba$ in
% this grammar.

\paragraph{Regular Expressions}
Regular expressions can be encoded as types generated by base types,
$\oplus$, and $\otimes$, and the Kleene star $(\cdot)^{\ast}$.  For a
grammar $\A$, we can define the Kleene star $\A^{*}$ as a particular
\emph{inductive linear type} of linear lists, as shown in
\cref{fig:kleenestarinductive}. Here $\A^{*} : \L$ means we are defining
a \emph{linear} type. $\A^{*}$ has two constructors: $\nil$, which
builds a parse of type $\A^{*}$ from nothing; and $\cons$, which
linearly consumes a parse of $\A$ and a parse of $\A^{*}$ and builds a
parse of $\A^{*}$. This linear consumption is defined by the linear
function type $\lto$. The linear function type $\A \lto \B$ defines functions that
take in parses of $\A$ as input, \emph{consume} the input, and return a parse of
$\B$ as output. The arrow, $\uparrow$, wrapping these
constructors means that the constructors by themselves are not consumed
upon usage, and so are \emph{non-linear} values themselves. That is, the
names $\nil$ and $\cons$ are function symbols that may be reused as many times
as we wish.

\begin{figure}
\begin{lstlisting}
data A(*@$^*$@*) : L where
  nil : (*@$\uparrow$@*)(A(*@$^*$@*))
  cons : (*@$\uparrow$@*)(A (*@$\lto$@*) A(*@$^*$@*) (*@$\lto$@*) A(*@$^*$@*))
\end{lstlisting}
% \begin{align*}
% \data &~A^{*} : L~\where\\
%       & \nil : \ltonl {A^{*}} \\
%       & \cons : \ltonl {(A \lto A^{*} \lto A^{*})}
% \end{align*}
\caption{Kleene Star as an inductive type}
\label{fig:kleenestarinductive}
\end{figure}

Through repeated application of the Kleene star constructors, \cref{fig:kleenestarderivation} gives a derivation that shows \stringquote{ab} matches the
regular expression $({\literal a}^{*} \otimes \literal b) \oplus \literal c$. The leaves of the proof tree
that mention the arrow $\uparrow$ describe a cast from a non-linear type to a linear type.
For instance, the premise of the leaf involving $\nil$
views $\nil : \ltonl {({\literal a}^{*})}$ as the name of a
constructor, and a constructor should be nonlinearly valued because we may call
it several time. However, the conclusion of this leaf views $\nil : {\literal a}^{*}$ as a
linear value. After we call the constructor it returns a value that may only be
used a single time. In this paper such casts are often made implicitly, but this
does not introduce any ambiguity.
Even though we use $\nil$ to refer to both a constructor and the
linear value produced by it
, we differentiate between the two views by inspecting
the type.

\steven{TODO fig 3 sizing}
\begin{figure}
\begin{mathpar}
\footnotesize
\inferrule
{
  \inferrule
  {
    \inferrule
    {
      \inferrule
      {
        \inferrule
        {
          \inferrule
          {~}
          {\cdot \vdash \cons : \ltonl {({\literal a} \lto {\literal a}^* \lto {\literal a}^*)}}
        }
        {\cdot \vdash \cons : {\literal a} \lto {\literal a}^* \lto {\literal a}^*}
        \\
        \inferrule
        {~}
        {a : {\literal a} \vdash a : {\literal a}}
      }
      {a : a \vdash \cons~a : {\literal a}^* \lto {\literal a}^*}
      \\
      \inferrule
      {
        \inferrule
        {~}
        {\cdot \vdash \nil : \ltonl {({\literal a}^*)}}
      }
      {\cdot \vdash \nil : {\literal a}^*}
    }
    {a : {\literal a} \vdash \cons~a~\nil : {\literal a}^{*}}
    \\
    \inferrule
    {~}
    {b : {\literal b} \vdash b : {\literal b}}
  }
  {a : {\literal a} , b : {\literal b} \vdash (\cons~a~\nil) \otimes b : {\literal a}^{*} \otimes {\literal b}}
}
{a : \literal a , b : \literal b \vdash \texttt{g} := \inl ((\cons~a~\nil)
  \otimes b) : ({\literal a}^{*} \otimes \literal b) \oplus \literal c}
\end{mathpar}
\begin{lstlisting}
g : (*@$\uparrow$@*)(('a' (*@$\otimes$@*) 'b') (*@$\lto$@*) ('a'(*@$^*$@*) (*@$\otimes$@*) 'b') (*@$\oplus$@*) 'c')
g (a , b) = inl (cons a' nil (*@$\otimes$@*) b')
\end{lstlisting}
\caption{\stringquote{ab} matches
  $({\literal a}^{*} \otimes \literal b) \oplus \literal c$}
\label{fig:kleenestarderivation}
\end{figure}

We may also have derivations where the term in context is not simply a
string of base types. In \cref{fig:kleeneabstractproof} we show that every parse
of the grammar $(A \otimes A)^{*}$ induces a parse of $A^{*}$ for an
arbitrary grammar $A$. The context $(A \otimes A)^{*}$ does not correspond directly to a string, so it is
not quite appropriate
to treat linear terms here as parse \textit{trees}.
The
context $a : (A \otimes A)^{*}$ does not contain concrete data to be parsed; rather, there may be many choices of string underlying the parse tree captured
by the variable $a$. Thus, the term $\texttt{h}$ from
\cref{fig:kleeneabstractproof} is not a parse of a string, and
it is more
appropriate to think of it as a parse \textit{transformer} --- a function from
parses of $(\A \otimes \A)^{*}$ to parse of $\A^{*}$.

We define $\texttt{h}$ by recursion on terms of type $(\A \otimes \A)^{*}$.
This recursion is expressed in the derivation tree by invoking the
elimination principle for Kleene star, written
as $\fold$. The parse transformer $\texttt{h}$ is more compactly presented in the pseudocode
of \cref{fig:kleeneabstractproof} by pattern matching on the input and making an
explicit recursive call in the body of its definition. When appropriate we will favor this code-like
presentation of terms in our calculus, as it is more readable and more
faithful to the implementation of \theoryabbv in Agda.

\begin{figure}
\begin{mathpar}
  \inferrule
  {
    \inferrule
    {~}
    {\cdot \vdash \nil : \A^*}
    \\
    \inferrule
    {
      \inferrule
      {
        \inferrule
        {~}
        {a_1 : \A, a_2 : \A , a_{*} : \A^* \vdash \cons (a_1 , \cons(a_2 , a_{*})) : \A^*}
      }
      {a' : \A \otimes \A , a_{*} : \A^* \vdash g := \letin {a_1 \otimes a_2} {a'} {\cons (a_1 , \cons(a_2 , a_{*})) : \A^*}}
    }
    {\cdot \vdash f := \lamblto a {\lamblto {a_{*}} {\app {\app g a} {a_{*}} }} : (\A \otimes \A) \lto \A^* \lto \A^*}
  }
  {a : (\A \otimes \A)^* \vdash \texttt{h} := \fold(\nil , f)(a) : \A^* }
\end{mathpar}
\begin{lstlisting}
h : (*@$\uparrow$@*)((A (*@$\otimes$@*) A)(*@$^*$@*) (*@$\lto$@*) A(*@$^*$@*))
h nil = nil
h (cons (a1 (*@$\otimes$@*) a2) as) = cons a1 (cons a2 (h as))
\end{lstlisting}
\caption{A parse transformer for abstract grammars}
\label{fig:kleeneabstractproof}
\end{figure}

\paragraph{Non-deterministic Finite Automata}
\newcommand{\s}{\texttt{s}}
\newcommand{\0}{\texttt{0}}
\newcommand{\1}{\texttt{1}}
\newcommand{\2}{\texttt{2}}
Regular expressions are a compact formalism for defining a formal grammar,
however simply defining a regular expression --- such as
$({\literal a}^{*} \otimes \literal b) \oplus \literal c$ --- does not give any insight into how to parse
it. Many parsing algorithms are based on a notion of automaton, and to verify
these algorithms we need to represent their associated automata in
\theoryabbv.

Finite automata are precisely the class of machines that recognize regular
expressions. \cref{fig:exampleNFA} shows a non-deterministic
finite automaton (NFA) for the regular expression
$({\literal a}^{*} \otimes \literal b) \oplus \literal c$.
\cref{fig:nfainductive} defines $\Trace$ an \emph{indexed} inductive linear type of
traces through this automaton. Defining an indexed inductive type can be thought
of like defining a family of mutually inductive types varying coherently over a
type of indexing data. That is, \cref{fig:nfainductive} defines the linear types
$\Trace~\0$,
$\Trace~\1$, and $\Trace~\2$ mutually inductively. The index $\s : \texttt{Fin 3}$ picks out which
state in the automaton a trace begins at --- where $\texttt{Fin 3}$ is the
finite type containing inhabitants $\{\0 , \1 , \2\}$.
There are three kinds of constructors for $\Trace$:
(1) those that terminate traces, (2) those that correspond to transitions
labeled by a character, and (3) those that correspond to transitions labeled by
the empty string $\epsilon$. The constructor $\texttt{stop}$ terminates a trace
in the accepting state $\2$. The constructors $\texttt{1to1}$, $\texttt{1to2}$,
$\texttt{0to2}$ each define a labeled transition through the NFA, and each of
these consumes a parse of the label's character and a trace beginning at the
destination of a transition to produce a trace beginning at the source of a
transition. The constructor $\texttt{0to1}$ behaves similarly, except its
transition is labeled with the empty string $\epsilon$. Therefore,
$\texttt{0to1}$ takes in a trace beginning a state $\1$ and returns a trace
beginning at state $\0$ corresponding to the same underlying string.

Just as \cref{fig:kleenestarderivation} shows the \stringquote{ab} matches
$({\literal a}^{*} \otimes \literal b) \oplus \literal c$, \cref{fig:nfaderivation} demonstrates that \stringquote{ab} is
accepted by the NFA representation of this regular expression. In particular, we
show that \stringquote{ab} matches the grammar $\Trace~0$ because $\0$ is the
initial state of the NFA.\ We may imagine defining parse transformers from
$({\literal a}^{*} \otimes \literal b) \oplus \literal c$ to $\Trace~\1$, and vice versa. The existence of
such maps between the two types would demonstrate that the regular expression
and the automaton capture the same language --- i.e.\ they recognize the same
set of strings. As we will see in LATER SECTION, we can define parse
transformers between regular expressions and their associated automaton.
Moreover, we prove that these parse transformers define a bijective
correspondence of parses over every input string. That is, in a strong sense the
parse of a regular expression and the traces of their associated automata
capture precisely the same data.

\begin{figure}
  \begin{tikzpicture}[node distance = 25mm ]
    \node[state, initial] (0) {$\0$};
    \node[state, below left of=0] (1) {$\1$};
    \node[state, right of=1, accepting] (2) {$\2$};

    \path[->] (0) edge[above] node{$\epsilon$} (1)
              (0) edge[right] node{$\stringquote{c}$} (2)
              (1) edge[loop left] node{$\stringquote{a}$} (1)
              (1) edge[below] node{$\stringquote{b}$} (2);
  \end{tikzpicture}
  \caption{NFA for $(\a^{*} \otimes \b) \oplus \c$}
  \label{fig:exampleNFA}
\end{figure}

\begin{figure}
\begin{lstlisting}
data Trace : (s : Fin 3) (*@$\to$@*) L where
  stop : (*@$\uparrow$@*)(Trace 2)
  1to1 : (*@$\uparrow$@*)('a' (*@$\lto$@*) Trace 1 (*@$\lto$@*) Trace 1)
  1to2 : (*@$\uparrow$@*)('b' (*@$\lto$@*) Trace 2 (*@$\lto$@*) Trace 1)
  0to2 : (*@$\uparrow$@*)('c' (*@$\lto$@*) Trace 2 (*@$\lto$@*) Trace 0)
  0to1 : (*@$\uparrow$@*)(Trace 1 (*@$\lto$@*) Trace 0)
\end{lstlisting}
\caption{$N$ as an indexed inductive type}
\label{fig:nfainductive}
\end{figure}

\begin{figure}
\begin{lstlisting}
k : (*@$\uparrow$@*)(('a' (*@$\otimes$@*) 'b') (*@$\lto$@*)  Trace 0)
k (a , b) = 0to1 (1to1 a (1to2 b stop))
\end{lstlisting}
\caption{\stringquote{ab} is accepted by the NFA for
  $({\literal a}^{*} \otimes {\literal b }) \oplus \literal c$}
\label{fig:nfaderivation}
\end{figure}

%% 3. Third example: NFA traces to show the indexed inductive types

\section{Syntax and Typing for \theoryname}
\label{sec:tt}

%% OUTLINE
%%
%% 1. Overview of syntax: what are the judgments
%% 2. Cover the linear types except for inductives
%% 3. Cover non-linear types, mainly focus on the universes
%% 4. Cover indexed inductive linear types
%% 5. Cover the "axioms" we add

The syntax of
\theoryabbv is based on the dependent linear-non-linear (\lnld)
calculus from \cite{krishnaswami_integrating_2015}. As in \lnld, \theoryabbv includes both
non-linear dependent types, as well as linear types, which are allowed
to depend on the non-linear types, but not on other linear types.
%
The main point of departure from \lnld's design in that the linear
typing is \emph{non-commutative}, i.e., that exchange is not an
admissible structural rule. Furthermore, we add a general-purpose
indexed inductive linear type connective, as well as an
\emph{equalizer} type, which we will show allows us to perform
inductive proofs of equalities between linear terms.
%
Finally, we add base types and axioms to the system specifically to
model formal grammars and parsing.

The formation rules for the judgments of \theoryabbv are shown in
\Cref{fig:formation}. $\Gamma$ stand for non-linear contexts, $X,Y,Z$
stand for non-linear types, $M,N$ stand for non-linear terms, these
act as in an ordinary dependent type theory. $\Delta$ stands for
linear contexts, $A,B,C$ for linear types and $e,f,g$ for linear
terms. These contexts, types and terms are allowed to depend on an
ambient non-linear context $\Gamma$, but note that linear types $A$
cannot depend on any \emph{linear} variables in $\Delta$. We include
definitional equality judgments for both kinds of type and term
judgments as well. Additionally, we have a judgment $\Gamma\vdash X
\isSmall$ which is used in the definition of universe types.

\begin{figure}
  \begin{mathpar}
    \inferrule{~}{\ctxwff \Gamma}

    \inferrule{\ctxwff \Gamma}{\ctxwffjdg \Gamma X}

    \inferrule{\ctxwffjdg \Gamma X}{\Gamma \vdash X \isSmall}

    \inferrule{\ctxwffjdg \Gamma X \and \ctxwffjdg \Gamma Y}{\ctxwffjdg \Gamma {X \equiv Y}}

    \inferrule{\ctxwffjdg \Gamma X}{\Gamma \vdash M : X}

    \inferrule{\Gamma \vdash M : X \and \Gamma \vdash N : X}{\Gamma \vdash M \equiv N : X}

    \inferrule{\ctxwff \Gamma}{\linctxwff \Gamma \Delta}

    \inferrule{\ctxwff \Gamma}{\Gamma \vdash A \isLinTy}

    \inferrule{\Gamma \vdash A \isLinTy \and \Gamma \vdash B \isLinTy}{\Gamma \vdash A \equiv B}

    \inferrule{\Gamma \vdash \Delta \isLinCtx \and \Gamma \vdash A \isLinTy}{\Gamma;\Delta\vdash e : A}

    \inferrule{\Gamma; \Delta \vdash e : A \and \Gamma;\Delta \vdash f : A}{\Gamma;\Delta\vdash e \equiv f : A}
  \end{mathpar}
  \caption{Formation Rules}
  \label{fig:formation}
\end{figure}

\subsection{Non-linear Typing}

The non-linear types of \theoryabbv include the standard dependent
types for $\Pi,\Sigma$, extensional equality, natural numbers,
booleans, unit and empty types. We include these rules in the
appendix. We present the other non-linear type constructions in
\Cref{fig:non-linear-stuff}. First, we include universe types $U$ of
small non-linear types and $L$ of linear types. These are defined as
``universes ala Coquand'' in that we define a judgment saying when a
non-linear type is small and define the universe to internalize
precisely this judgment. These universe types are needed so that we
can define types by recursion on natural numbers. Next, we include a
non-linear type $\ltonl A$ where $A$ is a linear type. The intuition
for this type is that its elements are the linear terms that are
``resource free'': its introduction rule says we can construct an
$\ltonl A$ when we have a linear term of type $A$ with no free linear
variables. This type is used extensively in our examples, playing a
similar role to the $!$ modality of ordinary linear logic.
\pedro{We should also reference the persistence modality $\square$
from separation logic}

%% \begin{figure}
%%   \[
%%   \begin{array}{rrcl}
%%     \textrm{non-linear contexts} & \Gamma & ::= & \cdot \pipe \Gamma,x:X \\
%%     \textrm{non-linear types} & X , Y , Z & ::= & \mathsf{Nat} \pipe
%%                                                   \mathsf{Bool} \pipe
%%                                                   \mathsf{Empty} \pipe
%%                                                   \mathsf{Unit} \pipe
%%                                                   \mathsf{SPF}\,X\pipe U \pipe L
%%                                                   \pipe \ltonl A \pipe \sum\limits_{x:X} Y
%%                                                   \pipe \prod\limits_{x:X} Y\pipe M =_A
%%                                                   N
%%   \end{array}
%%   \]
%%   \caption{Non-linear Types}
%% \end{figure}
\begin{figure}
  \begin{mathpar}
    \inferrule{}{\Gamma \vdash U \isTy}\and
    \inferrule
    {\Gamma \vdash M : U}
    {\Gamma \vdash \unquoteTy M \isTy}\and
    \inferrule
    {\Gamma \vdash X \isTy}
    {\Gamma \vdash \quoteTy X : U}\and
    \unquoteTy{\quoteTy{X}} \equiv X \and \inferrule{\Gamma \vdash M : U}{\Gamma\vdash \quoteTy{\unquoteTy{M}}\equiv M : U}

    \inferrule{}{\Gamma \vdash L \isTy}\and
    \inferrule
    {\Gamma \vdash M : L}
    {\Gamma \vdash \unquoteTy M \isLinTy}\and
    \inferrule
    {\Gamma \vdash A \isLinTy}
    {\Gamma \vdash \quoteTy A : L}\and
    \unquoteTy{\quoteTy{A}} \equiv A \and \inferrule{\Gamma \vdash M : U}{\Gamma\vdash \quoteTy{\unquoteTy{M}}\equiv M : U}

    \inferrule
    {\Gamma \vdash A \isLinTy}
    {\Gamma \vdash \ltonl { A } \isTy}
    \and
    \inferrule{\Gamma ; \cdot \vdash e : A}{\Gamma \vdash e : \ltonl A}\and
    \inferrule{\Gamma \vdash M : \ltonl A}{\Gamma ; \cdot \vdash M : A} \and
  \end{mathpar}
  \caption{Non-linear types (selection)}
  \label{fig:non-linear-types}
\end{figure}

\subsection{Linear Typing}

We give an overview of the linear types in \Cref{fig:linear-types} and
terms in \Cref{fig:linear-terms}. The equational theory for these
types is straightforward and included in the appendix. First, the
linear variable rule says that a linear variable can be used if it is
the \emph{only} variable in the context.

First, we cover the ``multiplicative'' connectives of non-commutative
linear logic. The linear unit and tensor product are standard for a
non-commutative linear logic: when we construct a linear unit we
cannot use any variables and when we construct a tensor product, the
two sides must use disjoint variables, and the variables the left side
of the product uses must be to the left in the context of the
variables used by the right side of the tensor product. Because we are
non-commutative, there are two function types: $A \lto B$ and $B \tol
A$. The difference between these is that the introduction rule for $A
\lto B$ adds a variable to the right side of the context, whereas the
introduction rule for $B \tol A$ adds a variable to the left side of
the context. In our experience, because by convention parsing
algorithms parse from left-to-right, we rarely need to use the $B \tol
A$ connective. One of the main uses of linear types is that, in
combination with the $\uparrow$ connective, we can abstract over
linear terms, in that the non-linear type $\ltonl{A \lto B}$ is a
non-linear type of linear terms with one linear input $A$ and output
$B$.
\pedro{Isn't this is a slightly incorrect desciption of $\ltonl{A \lto B}$?
How is the description above distinct form $A \lto B$?}

Next, we have the ``additive'' connectives. For these, we use the
non-linear types to define \emph{indexed} versions of the additive
disjunction $\oplus$ and additive conjunction $\&$ of linear
logic. These are like a linear version of the $\Sigma$ and $\Pi$
connectives of ordinary dependent type theory. The indexed $\&$ is
defined by a $\lambda$ and eliminated using projection. The indexed
$\oplus$ is like a \emph{weak} $\Sigma$ type: it has an injection
introduction rule $\sigma$, but its elimination rule is given by
\emph{pattern matching} rather than projections. Note that we can
define the more typical nullary and binary versions of these
connectives by using indexing over the empty and boolean type
respectively. We will freely use $\bot$ to refer to this empty
disjunction and $\top$ to refer to the empty conjunction. Lastly, we
include a type $\equalizer {a}{f}{g}$ that we call the
\emph{equalizer} of $f$ and $g$. We think of this type as the
``subset'' of $A$ that satisfy the equation $f\,a\equiv g\,a$. If
these were non-linear types, such a type could be constructed using
the equality type as $\sum_{x:X} f\,x=_Y g\,x$, but this construction
can't be used because it uses a \emph{dependent} equality type. We
note that the equalizer type is not used in defining any of our
parsers or formal grammars, but instead will be used shortly to derive
an induction principle for proving equalities of functions whose
domain is an inductive type.

\begin{figure}
  \begin{mathpar}
    \inferrule{}{\Gamma \vdash I \isLinTy}\and
    \inferrule{\Gamma \vdash A \isLinTy \and \Gamma \vdash B \isLinTy}{\Gamma \vdash A \otimes B \isLinTy}\and
    \inferrule{\Gamma \vdash A \isLinTy \and \Gamma \vdash B \isLinTy}{\Gamma \vdash A \lto B \isLinTy}\and
    \inferrule{\Gamma \vdash A \isLinTy \and \Gamma \vdash B \isLinTy}{\Gamma \vdash A \tol B \isLinTy}\and
    \inferrule{\Gamma,x:X \vdash A \isLinTy}{\Gamma \vdash \bigoplus\limits_{x:X} A}\and
    \inferrule{\Gamma,x:X \vdash A \isLinTy}{\Gamma \vdash \bigamp\limits_{x:X} A}\and
    \inferrule
    {\Gamma \vdash f : \ltonl {(A \lto B)} \and \Gamma \vdash g : \ltonl { (A \lto B) }}
    {\Gamma \vdash \equalizer {a}{f}{g} \isLinTy}
  \end{mathpar}
  \caption{Linear Types (selection)}
\end{figure}

\begin{figure}
  \begin{mathpar}
    \inferrule{~}{\Gamma ; a : A \vdash a : A}
    \and
    \inferrule{\Gamma ; \Delta \vdash e : B \\ \linctxwffjdg \Gamma {A \equiv B}}{\Gamma ; \Delta \vdash e : A}
    %
    \\
    %
    \inferrule{~}{\Gamma ; \cdot \vdash () : I}
    \and
    \inferrule{\Gamma ; \Delta \vdash e : I \\ \Gamma ; \Delta_1',\Delta_2' \vdash e' : C}{\Gamma ; \Delta_1',\Delta,\Delta_2' \vdash \letin {()} e {e'} : C}
    %
    \\
    %
    \inferrule{\Gamma ; \Delta \vdash e : A \\ \Gamma ; \Delta' \vdash e' : B}{\Gamma ; \Delta, \Delta' \vdash e \otimes e' : A \otimes B}
    %
    \and
    %
    \inferrule{\Gamma ; \Delta \vdash e : A \otimes B \\ \Gamma ; \Delta'_1, a : A, b : B, \Delta'_2 \vdash e' : C}{\Gamma ;  \Delta_1', \Delta, \Delta'_2 \vdash \letin {a \otimes b} e {e'} : C}
    \\
    %
    \inferrule{\Gamma ; \Delta , a : A \vdash e : B}{\Gamma ; \Delta \vdash \lamblto a e : A\lto B}
    \and
    \inferrule{\Gamma ; \Delta' \vdash e' : A \\ \Gamma ; \Delta \vdash e : A \lto B}{\Gamma ; \Delta, \Delta' \vdash \applto {e'} {e} : B}
    \\
    %
    \inferrule{\Gamma, x : X ; \Delta  \vdash e : A}
              {\Gamma ; \Delta \vdash \dlamb x e : \LinPiTy x X A}
    \and
    \inferrule{\Gamma ; \Delta \vdash e : \LinPiTy x X A \\ \Gamma \vdash M : X}{\Gamma ; \Delta \vdash e\,.\pi\,M : \subst A {e'} x}
    %
    \\
    %
    \inferrule{\Gamma \vdash M : X \quad \Gamma ; \Delta \vdash e : \subst A e x}{\Gamma ; \Delta \vdash \sigma\,M\,e : \bigoplus\limits_{x:X} A}
    %
    \and
    %
    \inferrule{\Gamma ; \Delta \vdash e : \bigoplus\limits_{x:X} A \quad \Gamma, x : X ; \Delta'_1, a : A, \Delta'_2 \vdash e' : C}{\Gamma; \Delta'_1, \Delta, \Delta'_2 \vdash \letin {\sigma\,x\,a} e {e'}: C}
    %
    \and
    %
    \inferrule
    { \Gamma ; \Delta \vdash e : A \\
      \Gamma ; \Delta \vdash \applto {f}{e} \equiv \applto {g}{e}}
    {\Gamma ; \Delta \vdash \equalizerin{e} : \equalizer {e}{f}{g}}
    %
    \and
    %
    \inferrule
    {\Gamma ; \Delta \vdash e : \equalizer{a}{f}{g}}
    {\Gamma ; \Delta \vdash \equalizerpi {e} : A}
  \end{mathpar}
  \caption{Linear terms (selection)}
  \label{fig:linsyntax}
\end{figure}
\pedro{Fixed tensor elimination rule (added output type)}
\pedro{Perhaps we should modify the equalizer introduction rule.
  As it stands, readers might be confused by the same symbol $e$
  being used everywhere.}

\subsection{Indexed Inductive Linear Types}

Next, we introduce one of the most complex and important linear
type constructors of our development, \emph{indexed inductive linear
types} $\mu F$. To encode these, we first add a non-linear type
$\SPF\,X$ of \emph{strictly-positive functorial} linear type
expressions indexed by a non-linear type $X$. We give the rules for
strictly-positive functors in \Cref{fig:spf}. First, we add operations
for a recursive reference $\Var~x$ the constant functor, $K$, tensor
products and additive conjunction and disjunction. As an example,
\max{TODO: make an example}.  Next we add an operation $\el(F)$ which
interprets our syntactic codes as an actual type constructor, and
define it for each of the operations given above. Lastly, we give a
\emph{functorial action} $\map(F)$ for each of the codes.\max{TODO:
  find a citation for this kind of approach. Container papers work
  this way for instance. I think there is a linear container paper
  that uses this approach}

With this notion of strictly-positive functor so defined, in
\Cref{fig:iilt} we next add a linear type constructor $\mu F$, which
defines a family of mutually-recursive inductive types for the family
of type expressions $F$. The introduction rule for this is
$\roll$, which constructs an element of $\mu F\,x$ from its
one-level unfolding. The elimination principle is defined by a
$\fold$ operation. The equations ensure that this makes the
family an \emph{initial algebra} for the type expressions. That is the
$\beta$ rule says that a $\mathsf{fold}$ applied to a $\mathsf{roll}$
is equivalent to $\mathsf{map}$ping the $\mathsf{fold}$ over all the
sub-expressions, which means that $\mathsf{fold}$ is an algebra
homomorphism from $\mathsf{roll}$ to $\mathsf{f}$. Then the $\eta$
rule says that $\mathsf{fold}$ is the \emph{unique} such homomorphism.

This definition as an initial algebra is well-understood semantically
but the $\eta$ principle in particular is somewhat cumbersome to use
directly in proofs. In dependent type theory, we would have a
dependent \emph{elimination} principle, which justifies
\emph{inductive proofs}. However since linear types do not support
dependency on linear types, we cannot directly adapt this
approach. However, if we are trying to prove that two morphisms out of
a mutually recursive type are equal, we can use the \emph{equalizer}
type to prove their equality by induction.\max{TODO: show how this is
  done.}

\begin{figure}
  \begin{align*}
    \mathsf{Var} &: X \to \SPF\,X\\
    \mathsf{K} &: L \to \SPF\,X\\
    \mathsf{\bigoplus} &: (Y \to \SPF\,X) \to \SPF\,X\\
    \mathsf{\bigamp} &: (Y \to \SPF\,X) \to \SPF\,X\\
    \mathsf{\otimes} &: \SPF\,X \to \SPF\,X \to \SPF\,X\\
    \el &: \prod_{X:U}\SPF\,X \to (X \to L) \to L\\
    \map &: \prod_{X:U}\prod_{F : \SPF\,X}\prod_{A,B:X\to L}(\prod_{x:X}\uparrow(A\,x\lto B\,x)) \to \uparrow(\el(F)(A) \lto \el(F)(B))\\\\
    \el(\Var\,M) B &= B\,M\\
    \el(\mathsf{K}\,A) B &= A\\
    \el(\bigoplus A) B &= \bigoplus_{y:Y}\el(A y)B\\
    \el(\bigamp A) B &= \bigamp_{y:Y}\el(A y)B\\
    \el(A \otimes A') B &= \el(A)B \otimes \el(A')B\\\\
    \map(\Var\,M)\,f &= f\,M\\
    \map(\K\,A)\,f &= \lambda a. a\\
    \map(\bigoplus A)\,f &= \lambda a. \letin{\oplusinj y {a_y}}{a}{\oplusinj y {\map(A\,y)\,f\,a_y}}\\
    \map(\bigwith\,A)\,f &= \lambda a. \withlamb{y}{\map(A\,y)\,f\,(\withprj y a)}\\
    \map(A \otimes A')\,f &= \lambda b. \letin{(a,a')}{b}{(\map(A)\,f\,a,\map(A')\,f\,a')}
  \end{align*}
  \begin{mathpar}
  \end{mathpar}
  \caption{Strictly Positive Functors}
  \label{fig:spf}
\end{figure}

\begin{figure}
  \begin{mathpar}
    \mathsf{roll} : \prod_{X:U}\prod_{F:X \to \SPF\,X}\prod_{x:X}\ltonl{(\el(F\,x)(\mu\,F))}

    %% \inferrule*[right=IndIntro]
    %% {\Gamma \vdash A : X \to \textrm{SPF}\,X\and
    %%  \Gamma \vdash M : X\and
    %%  \Gamma; \Delta \vdash e : \textrm{el}(A\,M) (\mu\,A)}
    %% {\Gamma; \Delta \vdash \mathsf{roll}\, e : \mu\,A\, M}

    \mathsf{fold} : \prod_{X:U}\prod_{F:X\to\SPF\,X}\prod_{A:X \to L}
    (\prod_{x:X}\ltonl{(\el(F\,x)(A) \lto A\,x)})
    \to \prod_{x:X}\ltonl{(\mu F\,x \lto A\,x)}

    \inferrule*[right=Ind$\beta$]
    {}
    {\fold\,F\,f\,x\,(\roll e) \equiv f\,x\,(\map(F\,x)\,(\fold\,F\,f))}

    \inferrule*[right=Ind$\eta$]
    {x:X;a:\el(F\,x)(\mu F) \vdash e'\,x\,(\roll a) \equiv f\,x\,(\map(F\,x)\,e')
    }
    {\fold\,F\,f\equiv e'}
    %% {\Gamma,x:X; a:\mu\,A\,x \vdash e' : B\,x
    %%   \and \Gamma;\Delta \vdash e : \mu A M\\\\
    %%   \Gamma,x:X; a':\el(A x)(\mu A) \vdash
    %%   e'[x/x][\roll(a')/a] = f[x/x][\map(A\,x)\,(y.a.e'[y/x][a/a])\,a/b] : B\,x
    %% }
    %% {\Gamma; \Delta \vdash \fold(A)(x.b.f)(M,e) = e'[M/x][e/a] : B\,M}
  \end{mathpar}
  \caption{Indexed Inductive Linear Types}
  \label{fig:iilt}
\end{figure}

\subsection{Grammar-specific Additions}

So far, our calculus is a generic combination of dependent types with
non-commutative linear types. In order to carry out formal grammar
theory and define parsers, however, we add a few carefully chosen base
types, terms and axioms that allow us to reason about formal grammars.
\max{TODO: more prose, also are these all of our axioms?}
\begin{enumerate}
\item (Characters) For each character of our alphabet $c \in \Sigma$, we add a linear type $c \isLinTy$.
\item (Strings are total) The function $\ltonl{(\String \lto \top)}$ is an isomorphism.\max{for this need to define $\String$, and isomorphism, and possibly we should define $\top$ in more detail earlier? idk}
\item (Distributivity of $\oplus,\&$) The function  $\ltonl{(\bigoplus\limits_{f : \prod_{x:X} Y(x)}\bigamp\limits_{x:X}A\,x\,(f\,x) \lto \bigamp\limits_{x:X}\bigoplus\limits_{y:Y(x)}A\,x\,y)}$ is an isomorphism.
\item ($\ltonl$ preserves $\oplus$) The function $\sum_{x:X}\ltonl{(A\,x)} \to \ltonl{(\bigoplus\limits_{x:X}A\,x)}$ is an isomorphism.
\end{enumerate}

% \steven{TODO the intro to this section is to be condensed based on what is
%   included in the examples section as motivation and intuition building}

% Omission of the structural rules of a deductive system, such as in
% linear logic \cite{GIRARD19871}, offers precise control over how a value is used
% in a derivation. Linear logic omits the weakening and contraction rules
% to ensure that every value in context is used exactly once. This control enables
% \emph{resource-sensitive} reasoning, where we may treat a resource as
% \emph{consumed} after usage. This viewpoint is amenable to parsing applications,
% as we may treat characters of a string as finite resources that are consumed at
% parse-time. That is, when reading a string, the occurrence of any character is read once. Freely duplicating or dropping characters from a string changes the meaning of that string. \max{what about a backtracking parser?}. One may then envision a linear type system where the types comprise
% formal grammars generated over some alphabet $\Sigma$, and the type constructors
% correspond precisely to inductive constructions on grammars --- such as
% conjunction, disjunction, concatenation, etc.

% Programming in a purely linear term language is limiting due to the variable
% usage restrictions. Code in such a language can
% become unnecessarily verbose when ensuring the linearity invariant.
% To alleviate this
% concern, in 1995 Benton et al.\ proposed an alternative categorical
% semantics of linear logic based on adjoint interactions between linear
% and non-linear logics \cite{bentonMixedLinearNonlinear1995} ---
% appropriately referred to as a \emph{linear-non-linear} (LNL)
% system. This work is simply typed, so the boundary between linear and
% non-linear subtheories is entirely characterized via a monoidal
% adjunction between linear terms and non-linear terms.

% \steven{Describe this monoidal adjunction with $\ltonl A$ and the derivable
%   version of $F$}

% Inside of an LNL system, linearity may be thought of as an option that users can
% choose to deliberately invoke at deliberate points in their developments in an
% otherwise intuitionistic system. However, if we are wishing to treat parsers as
% linear terms over input strings, the non-linear fragment of an LNL theory does
% not really assist in the development of parsers. It is instead the case that
% parsers may benefit from a \emph{dependence} on non-linear terms.
% Through the approach described by Krishnaswami et al.\ in
% \cite{krishnaswami_integrating_2015},
% we may define a restricted form of dependent types.

% In particular, dependence
% on linear terms
% is disallowed; however, through dependence of a linear term on a non-linear
% index, we can recover the definition of Aho's indexed grammars \cite{AhoIndexed}
% internal to \theoryabbv. Although, we can define strictly more things, as Aho's
% indexed grammars provide restricted access to the index whereas we can do
% whatever with it.

% \steven{elaborate on indexed grammars in a more articulate way in the last sentence}
% \steven{Put the aho stuff in the subtheory section}
%

% \paragraph{Regular Expressions}
% We may realize regular expression as a subtheory
% of our language by restricting the type constructors to the
% linear unit, characters, $\otimes$, $\oplus$, and Kleene star. Classically, the
% languages recognized by these are exactly the regular languages --- the lowest
% on the Chomsky hierarchy.

% \paragraph{$\mu$-Regular Expressions} Similarly, we may restrict the connectives
% to be linear unit, characters, $\otimes$, $\oplus$, and arbitrary recursion via
% left-fixed point rather than only Kleene star.
% Instead of regular expressions these correspond to Lei{\ss}'s $\mu$-regular
% expressions, which are known to be equivalent to context-free grammars
% \cite{leiss,krishnaswami_typed_2019}.

% \paragraph{Beyond Context-Free Grammars} The previous two subtheories induce as semantics
% regular grammars and context-free grammars, respectively; however, by including
% the LNL dependent types we may actually express the entirety of the Chomsky
% hierarchy. Through use of the LNL dependent types, we may encode indexed
% grammars, which are known to be properly between context-free and
% context-sensitive grammars within the Chomsky hierarchy \cite{AhoIndexed}. We
% will further see in \cref{subsubsec:tm} how to use dependence to encode Turing machines internal to our
% calculus, which of course induces a semantics of unrestricted grammars.

% \paragraph{Indexed Grammars}
% The LNL dependent types can be used to encode Aho's indexed grammars
% \cite{AhoIndexed} internal to \theoryabbv. Although, we can define strictly more
% things with this non-linear index, as we have unrestricted access to the index
% whereas Aho's grammars can only manipulate it in a particular way.

% \steven{Elaborate on these, and note that they are bw context free and context
%   sensitive on the Chomsky hierarchy}

\section{Formal Grammar Theory in \theoryname}
\label{sec:applications}
This section explores the applications of \theoryabbv to conducting formal
grammar theory. We demonstrate that several classical notions and constructions
integral to the theory of formal languages are faithfully represented
in \theoryname. By encoding well-established formal grammar concepts, we ensure
that our framework remains grounded in the
foundational principles of formal language theory while opening the door to
compositional formal verification of parsers.

\paragraph{Equivalence}
In \cite{chom1963}, Chomsky differentiates between the weak generative capacity
of a formal grammar --- the set of strings accepted by it --- and its strong
generative capacity --- the set of ``structural descriptions'', or parse trees,
that it generates. Two grammars are said to be \emph{weakly equivalent} if they
agree in their weak generative capacity, and likewise \emph{strongly equivalent}
if they agree in their strong generative capacity. We aim to recover each of
these notions of equivalence in \theoryabbv.

\begin{definition}
  \label{def:weakequiv}
  Grammars $\A$ and $\B$ are \emph{weakly equivalent} if there exists parse
  transformers $\f : \ltonl{(\A \lto \B)}$ and $\g : \ltonl {(\B \lto \A)}$.
\end{definition}
\pedro{Up to this point, we have mentioned a couple of times parse transformers,
  but we have never defined what they are. In particular, why does the definition
  require the $\uparrow$ modality. I'm assuming that it's to properly deal with
  linear variables, but these things are not obvious at first.}

\newcommand{\lang}{\texttt{Lang}}
A parse transformer $\f : \ltonl{(\A \lto \B)}$ takes as input parses of $\A$ for string
$\w$ and returns a parse of $\B$ for $\w$. Therefore, $\f$ demonstrates an
inclusion of languages $\lang(\A) \subset \lang(\B)$. Similarly, the existence of
$\g$ shows that $\lang(\B) \subset \lang(\A)$, and thus $\lang(\A) = \lang(\B)$.

\steven{Need better notation for language that is cohesive throughout the paper}

\begin{definition}
  \label{def:strongequiv}
  Grammars $\A$ and $\B$ are \emph{strongly equivalent} if there exists parse
  transformers $\f : \ltonl{(\A \lto \B)}$ and $\g : \ltonl {(\B \lto \A)}$ such
  that $\f \circ \g \equiv \ident$ and $\g \circ \f \equiv \ident$.
\end{definition}

For a string $\w$, \cref{def:strongequiv} induces a bijection between the
$\A$-parses of $\w$ and the $\B$-parses of $\w$, thus demonstrating that $\A$
and $\B$ agree in their strong generative capacity.

\paragraph{Ambiguity}
A formal grammar $\A$ is ambiguous if there is a string $\w$ for which $\A$ has
several parse trees of $\w$.
For example, $\a \oplus \a$ is ambiguous because
$\stringquote{a}$ has two parse trees --- on either side of the disjunction. On
the other hand, a formal grammar is unambiguous when there is at most
one parse tree for any input string. These notions of ambiguity are entirely
semantic, so we seek to find a property expressible in the syntax of \theoryabbv
that reifies the semantic description of ambiguity.

\begin{definition}
  \label{def:unambig}
  A grammar $\A$ is \emph{unambiguous} if for every linear type $\B$,
  $\f : \ltonl{(\B \lto \A)}$, and $\g : \ltonl{(\B \lto \A)}$
  then $\f \equiv \g$
\end{definition}
\cref{def:unambig} can be read more intuitively as stating that $\A$ is
unambiguous if there is at most one way to transform parses of any other grammar
$\B$ into parse of $\A$.

Note that this definition of unambiguity is preserved by strong equivalence, so a common proof technique for showing $\A$ is unambiguous is by proving $\A$
is strongly equivalent to an unambiguous grammar $\B$. In this style of proof, we
build the functions $\f : \ltonl{(\A \lto \B)}$ and
$\g : \ltonl{(\B \lto \A)}$, then we have the obligation to show
$\f \circ \g \equiv \ident$ and $\g \circ \f \equiv \ident$. However, because $\B$ is
unambiguous there is only one function of type $\ltonl{(\B \lto \B)}$, so
$\f \circ \g : \ltonl{(\B \lto \B)}$ and $\ident : \ltonl{(\B \lto \B)}$ are equal.
Thus, in this type of proof, we really only have the obligation of building each of $\f$
and $\g$ and demonstrating that $\g \circ \f \equiv \ident$.

% \paragraph{A Grammar of Strings}
% Define the character grammar $\charg$ to be the disjunction over characters in the alphabet $\Sigma$,
% $\charg = \LinSigTy {c} {\Sigma} c$. Then define a linear type of strings as a
% linear list of characters,
% $\stringg = \charg^{*}$.

% In the grammars model, we can build a
% strong equivalence between $\top$ and $\stringg$. hen building parsers, we may wish to enforce this
% syntactically by including $\top \cong \stringg$ as an axiom.

% Through inclusion of this axiom, $! : \stringg \to \top$ is an isomorphism, and
% thus also a monomorphism. That is, $\stringg$ is unambiguous.
% \max{we should have a section of the previous section which gives the axioms, this one and the distributivity axiom}

% \paragraph{Parseability}
% We say $A$ is \emph{parseable} if there exists a grammar $B$ and a map
% $\top \to A \oplus B$. Taking $\top \cong \stringg$, a grammar $A$ is parseable
% exactly when we may read an arbitrary string into a parse of $A$ or a $B$-valued
% refutation.

% \paragraph{Decidability}
% In the case that $A$ is parseable with respect to $\neg A$ --- i.e.\
% there's a map $\top \to A \oplus \neg A$ --- then call $A$ \emph{decidable}.


\subsection{Regular Expressions and Finite Automata}
%% 2. Regex, NFA, DFA
In this section, we make the equivalences
between regular expressions and finite automata precise by encoding two theorems
from classical language theory: (1) every regular expression is strongly equivalent
to the traces through some non-deterministic finite automaton (NFA); and, (2) for
every NFA, there is some deterministic finite automaton (DFA) such that the type
of traces through the NFA is weakly equivalent to the type of traces through the DFA.

\newcommand{\states}{\mathtt{states}}
\newcommand{\labelt}{\mathtt{label}}
\newcommand{\transitions}{\texttt{transitions}}
\newcommand{\epstransitions}{\epsilon\texttt{transitions}}
\newcommand{\isAcc}{\mathtt{isAcc}}
\newcommand{\init}{\mathtt{init}}
\newcommand{\src}{\mathtt{src}}
\newcommand{\dst}{\mathtt{dst}}
\newcommand{\parse}{\mathtt{parse}}
\newcommand{\print}{\mathtt{print}}
\newcommand{\epssrc}{\epsilon\mathtt{src}}
\newcommand{\epsdst}{\epsilon\mathtt{dst}}

\newcommand{\N}{\texttt{N}}

\paragraph{Regular Expressions} The theory of regular expressions emerges as the
subtheory of \theoryabbv generated by
$\otimes$, $\oplus$, and Kleene star over the types $\bot$, $I$, and characters drawn from $\Sigma$.

\paragraph{Nondeterministic Finite Automata} In \cref{sec:type-theory-examples}, we saw
one particular NFA corresponding to the regular expression
$({\literal a}^{*} \otimes \literal b) \oplus \literal c$. The definition of the linear type of traces
through this
NFA, as
presented in \cref{fig:nfainductive}, inspects the concrete structure of the
NFA.\ When generalizing from this example to an arbitrary NFA, we do not have the
luxury of inspecting a concrete definition. Therefore, we abstract over the
specification of an arbitrary NFA and define traces with respect to that
specification.

Let the data of an NFA $\N$ be given by the following nonlinear data,
\begin{itemize}
\item $\N.\states$ --- a finite set of states
\item $\N.\init : \N.\states$ --- the initial state
\item $\N.\isAcc : \N.\states \to \Bool$ --- a mapping marking which states are
        accepting
\item $\N.\transitions$ --- a finite set of labeled transitions
\item $\N.\labelt : \N.\transitions \to \Sigma$ --- a mapping from $N.\transitions$ to the character labeling the transition
\item $\N.\src , \N.\dst : \N.\transitions \to \N.\states$ --- mappings from
        $\N.\transitions$ to the source and target states, respectively, of a labeled transition
\item $\N.\epstransitions$ --- a finite set of $\varepsilon$-transitions
\item $\N.\epssrc, \N.\epsdst : \N.\epstransitions \to \N.\states$ --- mappings
        from $\N.\epstransitions$ to the source and target states of an $\epsilon$-transition
\end{itemize}

In \cref{fig:nfatrace} we define a linear type of traces
through an NFA $\N$. $\Trace_{\N}$ is an inductive type indexed by the first
state $\s : \N.\states$ of a trace and a Boolean $b : \Bool$ that marks if the
trace ends in an accepting state. Traces in $\N$ may be built through one of three constructors. We
may terminate a trace at an accepting state with the constructor $\nil$. If we
had a trace beginning at the destination state of a transition, then we may use
the $\cons$ constructor to linearly combine that trace with a parse of the label
of the transition to build a trace beginning at the source of the transition.
Finally, if we had a trace beginning at the destination of an
$\epsilon$-transition then we may use $\epscons$ to pull it back along the
$\epsilon$-transition and construct a trace beginning at the source of the
$\epsilon$-transition. As a shorthand, write $Parse_{N}$ for the accepting
traces out of $N.init$.

\begin{figure}
\begin{lstlisting}
data Trace(*@$_{\N}$@*) : (s : N.states) (b : Bool) (*@$\to$@*) L where
  nil : (*@$\uparrow$@*)(&[ s (*@$\in$@*) N.states ] Trace(*@$_{\N}$@*) s (N.isAcc s))
  cons : (*@$\uparrow$@*)(&[ t (*@$\in$@*) N.transitions ] &[ b (*@$\in$@*) Bool ]
           ('N.label t' (*@$\lto$@*) Trace(*@$_{\N}$@*) (N.dst t) b (*@$\lto$@*) Trace(*@$_{\N}$@*) (N.src t) b))
  (*@$\epsilon$@*)cons : (*@$\uparrow$@*)(&[ t (*@$\in$@*) N.(*@$\epsilon$@*)transitions ] &[ b (*@$\in$@*) Bool ]
           (Trace(*@$_{\N}$@*) (N.(*@$\epsilon$@*)dst t) b (*@$\lto$@*) Trace(*@$_{\N}$@*)(N.(*@$\epsilon$@*)src t) b))
\end{lstlisting}
\caption{Traces through an NFA $N$ as an indexed inductive type}
\label{fig:nfatrace}
\end{figure}

\paragraph{Deterministic Finite Automata}
\newcommand{\D}{\texttt{D}}
Just as the above defines an NFA, we provide a specification for an arbitrary
deterministic finite automaton.\ A DFA $D$ is
given by the following data,

\begin{itemize}
\item $\D.\states$ --- a finite set of states
\item $\D.\init : \D.\states$ --- the initial state
  \item $\D.\isAcc : \D.\states \to \Bool$ --- a mapping marking which states
        are accepting
\item $\D.\delta : \Sigma \to \D.\states \to \D.\states$ --- a deterministic transition function. Given a character $\c$ and a state
        $\s$, $\D.\delta~\c~\s$ is the state for which there is a transition $\s \overset{\stringquote{c}}{\to} (\D.\delta~\c~\s)$
\end{itemize}

\begin{figure}
\begin{lstlisting}
data Trace(*@$_{\D}$@*) : (s : D.states) (b : Bool) (*@$\to$@*) L where
  nil : (*@$\uparrow$@*)(&[ s (*@$\in$@*) D.states ] Trace(*@$_{\D}$@*) s D.isAcc s)
  cons : (*@$\uparrow$@*)(&[ c (*@$\in$@*) (*@$\Sigma$@*) ] &[ s (*@$\in$@*) D.states ] &[ b (*@$\in$@*) Bool ]
             ('c' (*@$\lto$@*) Trace(*@$_{\D}$@*) (D.(*@$\delta$@*) c s) b (*@$\lto$@*) Trace(*@$_{\D}$@*) s b))
\end{lstlisting}
\caption{Traces through the DFA $D$ as an indexed inductive type}
\label{fig:dfatrace}
\end{figure}

$\Trace_{\D}$, the linear type of traces through $\D$, is given in
\cref{fig:dfatrace}. These
traces may be terminated in an accepting state $\s$ with the $\nil$ constructor.
The $\cons$ constructor builds a trace out of state $\s$ by linearly combining
a parse of some character $\c$ with a
trace out of the state $\D.\delta~\c~\s$. The trace built with $\cons$ is
accepting if and only if the trace out of $D.\delta~c~s$ is accepting.

Because DFAs are deterministic, we are able to prove that their type of traces
are unambiguous and define a parser for them. In particular we show that
$\LinPiTy {s} {D.\states} {\LinSigTy {b} {\Bool} \Trace_{D}~s~b}$ is strongly
equivalent to
$\stringg$, which simultaneously constructs a parser and demonstrates that
traces through $D$ is unambiguous. Intuitively this strong equivalence says that we may
unambiguously build traces through $D$ beginning at any state $s$, and further
that from any trace through $D$ we can uniquely recover the string that built
the trace.

The first half of the equivalence,
$\parse_{D}$,
is defined by induction on strings in \cref{fig:dfaparser}. If this string is empty, then $\parse_{D}$
defines a linear function that terminates a trace at the input state $s$. If the
string is nonempty, then $\parse_{D}$ walks forward in $D$ from the input state $s$ by
the character at the head of the string.

\begin{figure}
\begin{lstlisting}
parse(*@$_D$@*) : (*@$\uparrow$@*)(String (*@$\lto$@*) &[ s (*@$\in$@*) D.states ] (*@$\oplus$@*)[ b (*@$\in$@*) Bool ] Trace(*@$_{\D}$@*) s b)
parse(*@$_D$@*) String.nil s = (*@$\sigma$@*) (D.isAcc s) (Trace(*@$_{\D}$@*).nil s)
parse(*@$_D$@*) (String.cons ((*@$\sigma$@*) c a) w) s =
   let (*@$\sigma$@*) b t = parse w (D.(*@$\delta$@*) c s) in
   (*@$\sigma$@*) b (Trace(*@$_{\D}$@*).cons c s b a t)
\end{lstlisting}
\caption{Parser for DFA traces}
\label{fig:dfaparser}
\end{figure}

The second half of the equivalence,
$\print_{D}$,
is given in \cref{fig:printdfa} by induction on traces. If the trace is defined
via $\nil$, then $\print_{D}$ returns the empty string. Otherwise, if the trace
is defined by $\cons$ then $\parse_{D}$ appends the character from the most
recent transition to the output
string and recurses.

\begin{figure}
\begin{lstlisting}
print(*@$_D$@*) : (s : D.states) (*@$\to$@*) (*@$\uparrow$@*)(((*@$\oplus$@*)[ b (*@$\in$@*) Bool ] Trace(*@$_{\D}$@*) s b) (*@$\lto$@*) String)
print(*@$_D$@*) s ((*@$\sigma$@*) b (Trace(*@$_{\D}$@*).nil .s)) = String.nil
print(*@$_D$@*) s ((*@$\sigma$@*) b (Trace(*@$_{\D}$@*).cons c (D.(*@$\delta$@*) c .s) b a trace)) =
  String.cons ((*@$\sigma$@*) c a) (print(*@$_D$@*) (D.(*@$\delta$@*) c s) ((*@$\sigma$@*) b trace))
\end{lstlisting}
\caption{Printer for DFA traces}
\label{fig:printdfa}
\end{figure}

\steven{TODO make the Agda logo clickable, and mention its meaning somewhere}
\newcommand{\Agda}{\agdalogo}
\begin{theorem}[Determinization, \Agda]
  \label{thm:determinization}
  Given an NFA $N$, there exists a DFA $D$ such that
  $Parse_{N}$ is weakly equivalent to $Parse_{D}$.
\end{theorem}
\newcommand{\X}{\texttt{X}}
\newcommand{\Y}{\texttt{Y}}
\newcommand{\x}{\texttt{x}}
\newcommand{\y}{\texttt{y}}
\begin{proof}
  We present a variant of Rabin and Scott's classic powerset construction \cite{rabinFiniteAutomataTheir1959}. First, we build
  the DFA $\D$. Define the
  states of $\D$ to be the $\mathbb{P}_{\varepsilon}(\N.\states)$ --- the type
  of $\epsilon$-closed\footnote{A subset of states $\X$ is $\epsilon$-closed if
    for every $\s \in \X$ and $\epsilon$-transition $\s \overset{\epsilon}{\to} \s'$ we have $\s' \in \X$.} subsets of
  $\N.\states$. A subset is accepting in $\D$ if it contains an accepting state
  from $\N$. Construct the initial state of $\D$ as the
  $\epsilon$-closure of $\N.\init$. Lastly define the transition function
  of $\D$ to send the subset $X$ under the character $\c$ to the
  $\epsilon$-closure of all the states reachable from $X$ via a transition
  labelled with the character $\c$.

  We demonstrate the weak equivalence between $Parse_{N}$ and
  $Parse_{D}$ by constructing parse transformers between the two
  grammars. To build the parse transformer
  $\ltonl{(Parse_{N} \lto Parse_{D})}$, we strengthen our inductive hypothesis
  and build the term
\newcommand{\NtoD}{\texttt{NtoD}}
\newcommand{\DtoN}{\texttt{DtoN}}
  \[
    \NtoD : \ltonl {\left(\Trace_{\N}~\s~\true \lto \bigwith\limits_{\X : \D.\states} \bigwith\limits_{\texttt{sInX} : \X \ni \s} \Trace_{\D}~\X~\true \right)}
  \]
  that maps a trace in $\N$ from an arbitrary state $\s$ to a trace in $\D$
  that may begin at any subset of states $\X$ that contains $\s$. $\NtoD$ may then be instantiated at
  $\s = \N.\init$ and $\X = D.\mathsf{init}$ to get the desired
  parse transformer.

\begin{figure}
\begin{lstlisting}
NtoD : (s : N.states) (*@$\to$@*)
       (*@$\uparrow$@*)(Trace(*@$_{\N}$@*) s true (*@$\lto$@*)
           &[ X (*@$\in$@*) D.states ] &[ sInX (*@$\in$@*) (X (*@$\ni$@*) s) ] Trace(*@$_{\D}$@*) X true)
NtoD .s (nil s) X sInX = nil s
NtoD .(N.src t) (N.cons trans .true a dstTrace) X sInX =
  D.cons (N.label t) X true a
    (((NtoD s dstTrace).(*@$\pi$@*) (D.(*@$\delta$@*) (N.label t) X)).(*@$\pi$@*) (N.dst t))
NtoD .(N.(*@$\epsilon$@*)src t) ((*@$\epsilon$@*)cons (*@$\epsilon$@*)trans .true (*@$\epsilon$@*)DstTrace) X sInX =
  ((NtoD s (*@$\epsilon$@*)DstTrace).(*@$\pi$@*) X).(*@$\pi$@*) (N.(*@$\epsilon$@*)dst t)
\end{lstlisting}
\caption{Parse transformer from NFA traces to Powerset DFA traces}
\label{fig:NtoD}
\end{figure}

Because the state $\s$ is now arbitrary, we can define $\NtoD$ via
recursion on the trace through the NFA. Code for this is given in \cref{fig:NtoD}.
If the NFA trace is terminating at the accepting state $\s$ and $\s \in \X$,
then we return a trace in $\D$ that terminates at
$\X$, which is accepting because it contains the accepting state $\s$.

If the NFA trace is defined via $\cons$ for a labeled transition
$t : N.transitions$ such that $N.src~t \in \X$, then
we first decompose the trace into a parse of $N.label~t$ and a
trace starting from $N.dst~t$. We then recursively call $\NtoD$
on this trace which returns a family of traces parameterized by $\epsilon$-closed
subsets containing $N.\dst~t$. We instantiate this family at the
set of states
$\D.\delta~(N.label~t)~X$. Because $\N.\dst~t$ is
reachable from $\N.\src~\texttt{trans}$ by a single transition labeled by
$N.label~t$ and $N.src~t \in X$ we have that
$N.dst~t \in \D.\delta~(N.label~t)~X$. Therefore
we can linearly combine the parse of the label and the trace built from the
recursive call to get a trace in $\D$ rooted at $X$.

Lastly, if the NFA trace is defined via $\epsilon\cons$ for an $\epsilon$-transition
$t : \N.\epstransitions$ such that
$\N.\epssrc~t \in \X$, then the construction is quite
similar. In this case we again recursively call $\NtoD$, except afterwards we
instantiate the family of grammars at the set $X$ itself. We may do this because
all states in $\D$ are taken to be $\epsilon$-closed, therefore
$\N.\epsdst~t \in \X$ because
$\N.\epssrc~t$ is in $X$. This concludes the construction of $\NtoD$.

To construct a term from DFA traces to NFA traces, we similarly strengthen our
induction hypothesis and build a parse transformer
\[
  \DtoN : \ltonl {\left( \Trace_{\D}~\X~\true \lto \bigoplus\limits_{s : \N.\mathsf{states}}\bigoplus\limits_{\texttt{sInX} : \X \ni \s} \Trace_{\N}~s~\true \right)}
\]

For a moment, let us turn our attention to the content behind the traces in $\D$.
Traces from subset $\X$ to subset $\Y$ describes the \emph{existence} of some trace
in the $\N$ from some $\x \in \X$ to some $\y \in \Y$. In general it is not clear how
to extract out a choice of trace in $\N$, which is what $\DtoN$ is asking of us.
However, we may further take the states and transitions of $\N$ to be \emph{ordered}.
This allows us to
induce an order on traces in $\D$ from $\X$ to $\Y$. From these orderings we can then
define the necessary choice functions on each type by taking the smallest with
respect to the ordering.

We define $\DtoN$ by recursion on traces through $\D$. First if the trace stops at
some accepting $\epsilon$-closed subset $\X$, then we may choose some state
$\s \in X$ such that $\s$ is accepting and terminate the trace in $\N$ at $\s$.

Next if the trace in $\D$ is defined via transition
$\X \overset{\stringquote{c}}{\to} (\D.\delta~\c~\X)$ then we may choose some state in
$\s \in \D.\delta~\c~\X$. Recursively, $\DtoN$ sends the trace rooted at
$\D.\delta~\c~\X$ to one rooted at $\s$. By definition of $\D.\delta~\c~\X$,
there exists some $\s' \in \X$ such that $\s'$ transitions to $\s$ by some number of
$\epsilon$-transitions followed by a transition labeled by $\c$.
We order all
such traces of this form and choose the smallest one, which is then linearly
combined with the parse of $\c$ to build an $\N$ trace. This concludes the
definition of $\DtoN$, thus showing that $Parse_{N}$ and $Parse_{D}$ are
weakly equivalent.
\end{proof}

\newcommand{\R}{\texttt{R}}
\begin{theorem}[Thompson's Construction, \Agda]
  Given a regular expression $\R$, there exists an NFA $\N$ such that $\R$ is
  strongly equivalent to $\Trace_{N}(\N.\init)$.
\end{theorem}

\begin{proof}
  \steven{TODO this proof needs to be rewritten. Likely through use of equalizers}
  \pedro{Since this proof and the determinization one are well-known and we have
    proved it in Agda, it's probably a good idea to include a much shorter proof sketch,
    or move it to the appendix, though this would feel a bit superfluous.}
We present a variant of Thompson's construction
\cite{thompsonProgrammingTechniquesRegular1968} in which we recurse on the
structure of $\R$ to construct a strongly equivalent NFAs. In this paper we will walk through
some of the cases of this proof. The remaining cases
follow
similar reasoning and the details may be found in the formalized artifact.

\begin{figure}
  \begin{tikzpicture}[node distance = 17mm ]
    \node[state, initial] (0) {$\0$};
    \node[state, right of=0, accepting] (1) {$\1$};

    \path[->] (0) edge[above] node{$\stringquote{c}$} (1);
  \end{tikzpicture}
  \caption{NFA $\N_{\c}$ for a single character $\c$}
  \label{fig:litnfa}
\end{figure}

First handle the case where $R$ is a literal character $\literal c$. The NFA $N_{c}$ pictured in
\cref{fig:litnfa} is strongly equivalent to the grammar $c$. To demonstrate this
strong equivalence, we build functions between $\literal c$ and $Parse_{N_{c}}$
and show they are mutually inverse.

\begin{figure}
\begin{lstlisting}
A : (s : Fin 2) (*@$\to$@*) L
A 0 = 'c'
A 1 = I

toN(*@$_c$@*) : (s : Fin 2) (*@$\to$@*) (*@$\uparrow$@*)(A s (*@$\lto$@*) Trace(*@$_{N_c}$@*) s true)
toN(*@$_c$@*) 0 c = cons t true c (nil 1)
toN(*@$_c$@*) 1 () = nil 1

fromN(*@$_c$@*) : (s : Fin 2) (*@$\to$@*) (*@$\uparrow$@*)(Trace(*@$_{N_c}$@*) s true (*@$\lto$@*) A s)
fromN(*@$_c$@*) 0 (cons t .true c tr) = let a (*@$\otimes$@*) () = c (*@$\otimes$@*) fromN(*@$_c$@*) 1 tr in a
fromN(*@$_c$@*) 1 (nil .1) = ()
\end{lstlisting}
\caption{Parse transformer out of $Trace_{N_{c}}~s~true$}
\label{fig:litNFAterms}
\end{figure}

In \cref{litNFATterms} we define a family of grammars $A$ and then prove that each $A~s$ is
equivalent to the accepting traces in $N_{c}$ out of state $s$. Then the desired
equivalence arises from taking $s = 0$.

When
defining $fromN_{c}$ we pattern match on the accepting traces beginning in each
state. Out of state $0$ the only accepting traces are built via $cons$ over the
single transition $t$, and out
of state $1$ the only accepting trace is built via $nil$.

Note that each $A~s$ is unambiguous, so $(fromN_{c}~s) \circ (toN_{c}~s)$ is equal to the
identity on each $A~s$. To complete the proof of strong equivalence it remains
to show that $(toN_{c}~s) \circ (fromN_{c}~s)$ is equal to the identity on
$Trace~s~true$, which we show by using the equalizer type.

\begin{figure}
\begin{lstlisting}
eqHom(*@$_{N_c}$@*) : (s : N(*@$_c$@*).state) (*@$\to$@*)
      (*@$\uparrow$@*)(Trace(*@$_{N_c}$@*) s true (*@$\lto$@*) {tr | (toN(*@$_c$@*) s (*@$\circ$@*) fromN(*@$_{c}$@*) s) tr = id tr})
\end{lstlisting}
\caption{Signature for the equalizer term}
\label{fig:litequalizer}
\end{figure}

The type signature for the family of functions, $eqHom_{N_{c}}$, into the equalizer type is given in
\cref{fig:litequalizer}. By defining an $Trace_{N_{c}}$-algebra homomorphism
into the equalizer, we enable a
proof strategy for equality by induction on traces through $N_{c}$. First consider the traces
out of state $1$. The only case to handle here is $nil 1$, and by definition
$toN_{c}~1$ and $fromN_{c}~1$ invert each other.
When defining $eqHom_{N_{c}}$ for traces out of state $0$, we must leverage an
inductive hypothesis,

\begin{align*}
  (toN_{c}~0 \circ fromN_{c}~0) (cons~t~true~c~tr)
  & \equiv (toN_{c}~0) (let~a~\otimes~()~=~c~\otimes~fromN_{c}~1~tr~in~a) \\
  & \equiv cons~t~true~(let~a~\otimes~()~=~c~\otimes~fromN_{c}~1~tr~in~a)~(nil~1) \\
  & \equiv cons~t~true~c~(nil~1) \\
  & \equiv cons~t~true~c~tr \\
\end{align*}

The first two equalities are by defintion. The thrid equality holds because
$\literal c$ is unambiguous and
$let~a~\otimes~()~=~c~\otimes~fromN_{c}~1~tr~in~a$ and $c$ have type
$\literal c$. The final equality holds because our inductive hypothesis states
that $nil~1 \equiv tr$. Thus, each $A s$ and $Trace_{N_{c}}~s~true$ are strongly
equivalent --- in particular, $\literal c$ and $Parse_{N_{c}}$.

\steven{Rewrite the disjunction case with equalizers}
Now consider the case where $R$ is the disjunction of two regular expressions
$R_{1}$ and $R_{2}$. Inductively, we may build NFAs $N_{1}$ and $N_{2}$ that are
strongly equivalent to $R_{1}$ and $R_{2}$, respectively. We define $N_{\oplus}$
as a new NFA built as the disjunction of $N_{1}$ and $N_{2}$.
$N_{\oplus}.\states$ comprises the states from $N_{1}$, the states from
$N_{2}$, and a new marked initial state. The accepting states are precisely
those that were accepting in the subautomata. The initial state has an
$\varepsilon$-transition to each of the initial states of the subautomata. Once in
$N_{1}.\init$, the downstream transitions are a copy of those in $N_{1}$,
likewise for $N_{2}.\init$.  That is, $N_{\oplus}$ contains a copy of each
$N_{1}$ and $N_{2}$, and we choose which copy to inhabit by deciding which
$\varepsilon$-transition out of the initial state to take. A schematic
representation of this automaton is given in \cref{fig:disjunctionnfa}.

\begin{figure}
  \begin{tikzpicture}[node distance = 17mm ]
    \node[state, initial] (init) {$\init$};
    \node[state, above right of=init] (init1) {$N_{1}.\init$};
    \node[state, below right of=init] (init2) {$N_{2}.\init$};
    \node[right of=init1] (dots1) {$\dots$};
    \node[right of=init2] (dots2) {$\dots$};
    \node[state, accepting, right of=dots1] (acc1) {$s_{1}$};
    \node[state, accepting, right of=dots2] (acc2) {$s_{2}$};

    \path[->] (init) edge[below] node{$\varepsilon$} (init1)
              (init) edge[above] node{$\varepsilon$} (init2)
              (init1) edge[above] node{} (dots1)
              (init2) edge[above] node{} (dots2)
              (dots1) edge[above] node{} (acc1)
              (dots2) edge[above] node{} (acc2);
  \end{tikzpicture}
  \caption{NFA $N_{\oplus}$ as a disjunction of $N_{1}$ and $N_{2}$}
  \label{fig:disjunctionnfa}
\end{figure}

Because $N_{1}$ and $N_{2}$ are each strongly equivalent to $R_{1}$ and $R_{2}$,
it suffices to show that
$Parse_{N_{1}} \oplus Parse_{N_{2}}$ is strongly
equivalent to $Parse_{N_{\oplus}}$.

Just like the case of a single character, when building the maps in either
direction, we strengthen the induction hypothesis by mapping out of traces
beginning with an arbitrary state then proceed by induction on traces.

First, define the
mapping out of $\Trace_{N_{\oplus}}(s)$. If $s$ is a state from $N_{1}$ then
send the trace to a trace in $N_{1}$ beginning at $s$. Likewise, if $s$ is from
$N_{2}$, then we map to a trace in $N_{2}$ beginning at $s$. If $s$ is
$N_{\oplus}.\init$, then we map the trace to a disjunction of traces
$\Trace_{N_{1}}(N_{1}.\init) \oplus \Trace_{N_{2}}(N_{2}.\init)$.

In the other direction, we need to map a disjunction of traces through $N_{1}$
or $N_{2}$ to a trace in $N_{\oplus}$. Without loss of generality, assume that
the trace came from $N_{1}$. That is, assume we currently have a term of type
$\Trace_{{N_{1}}}(N_{1}.\init)$ and we aim to construct a term of type
$\Trace_{N_{\oplus}}(N_{\oplus}.\init)$. We induct on the trace $N$ trace an map
it over to a trace beginning at the copy of $N_{1}.\init$ in $N_{\oplus}$, and
then we conclude by using the $\epscons$ constructor to turn this trace into one
beginning at $N_{\oplus}.\init$.

We show that these maps mutually invert each other by the same sort of argument
we used with the single character NFA.\ That is, we prove that the each map is a
homomorphism of algebras and then appeal to the initiality of each of
$\Trace_{N_{1}}$, $\Trace_{N_{2}}$, and $\Trace_{N_{\oplus}}$ to prove that the
relevant maps compose to the identity.

In this sketch, we elide the proof that these maps do indeed form homomorphisms.
Further, we are eliding the needed NFAs for the cases where the regular
expression is $I$, $\bot$, a Kleene star, or a concatenation. All of these
elided details may be found in the Agda formalization.
\end{proof}

\subsection{The Dyck Grammar}
%% 3. Dyck language
\newcommand{\balanced}{\mathtt{balanced}}
\newcommand{\Dyck}{\mathtt{Dyck}}
Moving up the Chomsky hierarchy from regular grammars, we arrive at context-free
grammars (CFGs). The smallest nontrivial CFG is the Dyck grammar
$\Dyck$ --- the grammar of balanced parentheses.

\pedro{Is it actually the smallest? How do you quantify size of nontrivial things?
  Might be better to just say ``an interesting nontrivial CFG [...]''}

$\Dyck$ is a grammar over the alphabet $\{ \stringquote{(}, \stringquote{)} \}$
and is defined in \cref{fig:dyckinductive}. The $\nil$ constructor shows that
the empty string is balanced, and the $bal$ constructor builds a balanced
parse by wrapping an already balanced parse in an additional set of parentheses
then following it with another balanced parse. We construct a parser for $\Dyck$ by building a deterministic automaton $M$ such
that $Parse_{M}$ is strongly equivalent to $\Dyck$.
\begin{figure}
\begin{lstlisting}
data Dyck : L where
  nil : (*@$\uparrow$@*) Dyck
  bal : (*@$\uparrow$@*)('(' (*@$\lto$@*) Dyck (*@$\lto$@*) ')' (*@$\lto$@*) Dyck (*@$\lto$@*) Dyck)
\end{lstlisting}
\caption{The Dyck grammar as an inductive linear type}
\label{fig:dyckinductive}
\end{figure}

\Cref{fig:DyckAutomaton} presents a sketch of the automaton $M$. The traditional
approach to parse grammars like $\Dyck$ is to employ a stack to count how many
open parentheses have been seen so far. We may think of
$M$ as performing the same task as a stack machine through slightly different means. Instead of storing
the number of open parentheses in a distinguished stack of memory, $M$ operates
by storing this data in the current state. $M$ is defined almost identically to how we defined DFAs, except we omit the
finiteness condition on the states of the automaton. This similarity allows
functions $parse_{M}$ and $print_{M}$ to be defined analogously to
\cref{fig:dfaparser,fig:printdfa}, and so
$\LinSigTy {s} {M.states} {\LinSigTy {b}{Bool} {Trace_{M}~s~b}}$ is likewise unambiguous.

\begin{theorem}[\Agda]
  \label{thm:dyck}
  $Dyck$ and $Parse_{M}$ are strongly equivalent.
\end{theorem}

We omit the details of \cref{thm:dyck} in this paper, and the ful proof is
mechanized in our Agda development. We can build a $Dyck$ parser by utilizing
this artifact, and the
runtime of the $Dyck$ parser is about 30 seconds for input that is 25,000
characters long.
\pedro{It might be a good idea to mention explicitly in the introduction
  that we are not claiming that our verified parsers are competitive with
  the state of the art.}

\begin{figure}
  \begin{tikzpicture}[node distance = 17mm ]
    \node[state, initial, accepting] (0) {$0$};
    \node[state, right of=0] (1) {$1$};
    \node[state, right of=1] (2) {$2$};
    \node[right of=2] (3) {$\dots$};
    \node[state, below of=0] (fail) {$fail$};

    \path[->] (0) edge[above, bend left] node{\stringquote{(}} (1)
              (1) edge[below, bend left] node{\stringquote{)}} (0)
              (1) edge[above, bend left] node{\stringquote{(}} (2)
              (2) edge[below, bend left] node{\stringquote{)}} (1)
              (2) edge[above, bend left] node{\stringquote{(}} (3)
              (3) edge[below, bend left] node{\stringquote{)}} (2)
              (fail) edge[loop right] node{\stringquote{(}, \stringquote{)}} (fail)
              (0) edge[left] node{\stringquote{)}} (fail);
  \end{tikzpicture}
  \caption{Automaton $M$ for the Dyck grammar}
  \label{fig:DyckAutomaton}
\end{figure}

% \subsection{LL(1)}
% %% 4. LL(1) parser

% \steven{TODO impute LL(1) example, mention that its weak equivalence is mechanized}

\subsection{Unrestricted Grammars}
\newcommand{\stringNL}{String'}
\steven{Need to better introduce nonlinear type of strings. We haven't talked
  about semantics nor have we included a nonlinear type of strings in the syntax}
Define $\stringNL$ to be the non-linear list type over the alphabet $\Sigma$ and
let $P : \stringNL \to U$ be a mapping from the set of strings to the universe of
nonlinear types. The function $P$ describes a family of nonlinear types
dependent on $\stringNL$.
We view the image of $P$ as a set of parse trees for each input string, and thus
$P$ carries the structure of a formal grammar. However, even though $P$ describes this grammatical data, it is not
a first-class linear type in \theoryabbv. To address this gap, below we define a linear type that captures the same
semantic grammar as $P$,

\[
Reify~P = \LinSigTyLimit {w} {\stringNL} {\LinSigTyLimit {x} {P~w} {~\lceil w \rceil}}
\]

\steven{Probably put this ceiling business around the discussion of the linear
  type of strings. Is this in section 3?}

where $\lceil \stringquote{} \rceil = I$ and
$\lceil c :: w \rceil = \literal c \otimes \lceil w \rceil$.

This reification operation on functions $\stringNL \to U$ is incredibly expressive,
as it allows to sidestep our linear typing connectives and utilize the whole of nonlinear
dependent type theory to define a grammar. This power allows us to encode
arbitrary grammars. In particular, given a
Turing machine $Tur$ one may
define a predicate $accepts : \stringNL \to U$ such that $accept~w$ is equal to $\top$ if
$T$ accepts $w$ and equal to $\bot$ otherwise. Then, $Reify~accepts$ is a linear
type that captures precisely the strings accepted by $Tur$. That is,
$Lang(Reify~accepts)$ is recursively enumerable --- the most general class of
languages in the Chomsky hierarchy.


\subsection{Semantic Actions}
So far, each construction has been primarily concerned with the parsing of a
string into a \emph{concrete} syntax tree. However, in a larger verified
developement --- such as a compiler --- a parser will often emit an
\emph{abstract} syntax tree that omits superflous syntactic details that aren't
needed for later passes.  The procedure that takes in a
concrete syntax tree and returns an abstract one is referred to as a
\emph{semantic action}.

Given a linear type $A$ and a nonlinear type $X$, we define a semantic action to
be a function $\ltonl {(A \lto \Delta X)}$, where
$\Delta X = \LinSigTy {x} {X} {\top}$. That is, a semantic action is an
elaboration from the concrete parses of $A$ into the semantic domain $X$. As
written, these actions are quite general. As discussed above, we may define
folds that are still grammatical in nature --- such as folding over a concrete
parse tree and forgetting parentheses that aren't needed in later compiler
passes. We may also define arbitrary computation over the structure of the parse
tree, such as calculating the maximum nesting depth of parentheses for $Dyck$,
returning the length the input string, or
even an interpreter that evaluates arithmetic expressions at parse-time.

\steven{What can I say about what we have in the Agda development here?}
\pedro{I'm assuming that this section is still incomplete, so I'll wait
for it to be complete in order to make more detailed comments.}

\section{Denotational Semantics}
\label{sec:semantics-and-metatheory}

To justify our assertion that \theoryabbv is a syntax for formal
grammars and parse transformers, we will now define a
\emph{denotational semantics} that makes this mathematically
precise. To interpret our types of codes we assume a universe of
\emph{small} sets.

\subsection{Formal Grammars and Parse Transformers}

First, we need to establish a precise definition of a formal grammar
that accounts for the possibility of ambiguity. The most common
definition of a formal grammar is as generative grammars, defined by
a set of non-terminals, with a specified start symbol and set of
production rules. We instead use a more abstract formulation that is
closer in spirit to the standard definition of a formal language:

\begin{definition}
  A \emph{formal language} $L$ is a function from strings to propositions.
  A \emph{formal grammar} $A$ is a function from strings to small sets.
\end{definition}

We think of the grammar $A$ as taking a string to the set of all parse
trees for that string. However since $A$ could be any function
whatsoever there is no requirement that an element of $A(w)$ be a
``tree'' in the usual sense. This definition provides a simple,
syntax-independent definition of a grammar that can be used for any
formalism: generative grammars, categorial grammars, or our own
type-theoretic grammars. Note that the definition of a formal grammar
is a generalization of the usual notion of formal language since a
proposition can be equivalently defined as a subset of a one-element
set. Then the difference between a formal grammar and a formal
language is that formal grammars can be \emph{ambiguous} in that there
can be more than one parse of the same string. Even for unambiguous
grammars, we care not just about \emph{whether} a string has a parse
tree, but \emph{which} parse tree it has, i.e., what the structure of
the element of $A(w)$ is.

The linear types of \theoryabbv can all be interpreted as formal
grammars, but what about the linear terms? These can be interpreted as
\emph{parse transformers}, which we now define.
\begin{definition}
  Let $L_1,L_2$ be formal languages. Then we say $L_1\subseteq L_2$ if
  for every string $w$, $L_1(w)\Rightarrow L_2(w)$.

  Let $A_1,A_2$ be formal grammars. Then a parse transformer $f$ from
  $A_1$ to $A_2$ is a function assigning to each string $w$ a function
  $f_w : A_1(w) \to A_2(w)$.
\end{definition}
\pa{Do we actually need a definition of formal language morphism?
In case we do, it should probably be its own definition.}
First, note that parse transformers generalize language inclusion: if
$A_1(w), A_2(w)$ are all subsets of a one-element set, then a parse
transformer is equivalent to showing that $A_1(w) \subseteq A_2(w)$.
In our denotational semantics, linear terms will be interpreted as
such parse transformers, and the notions of unambiguous grammar,
parseability, etc, introduced in \Cref{sec:concepts} can be verified
to correspond to their intended meanings under this interpretation.

Parse transformers can be composed: given two parse transformers $f$ and $g$,
their composition is defined pointwise, i.e. $(f\circ g)_w = f_w \circ g_w$.
Furthemore, given a formal grammar $A$, its identity transformer is $id_w =
id_{A(w)}$, where $id_{A(w)}$ is the identity function on the set $A(w)$. This
defines a \emph{category}.
\begin{definition}
  Define $\Grammar$ to be the category whose objects are formal
  grammars and morphisms are parse transformers.
\end{definition}
This category is equivalent to the slice category $\Set/\Sigma^*$ and
as such is very well-behaved. It is complete, co-complete, Cartesian
closed and carries a monoidal biclosed structure. We will use these
structures to model the linear types, terms and equalities in
\theoryabbv.
%% \begin{theorem}
%%   The category $\Grammar$ is complete, co-complete and cartesian
%%   closed. Furthermore, it is monoidal biclosed with the monoidal
%%   product given by
%%   \[ (A \otimes B) w = \{ (w_1,w_2,a,b) \pipe w_1w_2 = w \wedge a \in A w_1 \wedge b \in B w_2\} \]
%%   and unit by
%%   \[ I w = \{ () \pipe w = \varepsilon \} \]
%% \end{theorem}

%% \max{todo: need to decide if we need this section here or if it's in the intro already}.
% The linear typing judgment in our syntax takes on the following schematic form
% $\Gamma ; \Delta \vdash a : A$. First, $A$ represents a \emph{linear type} in
% our syntax. The intended semantics of these linear types are formal grammars.
% That is, the linear typing system is designed to syntactically reflect the
% behavior of formal grammars. For this reason, we may often interchangeably use
% the terms ``linear type'' and ``grammar''.

% The term $a$ is an inhabitant of type $A$, which is thought of as a parse tree of
% the grammar $A$. The core idea of this entire paper follows precisely from this
% single correspondence: grammars are types, and the inhabitants of these types
% are parse trees for the grammars.

% $\Gamma$ represents a non-linear context, while $\Delta$ represents a
% \emph{linear context} dependent on $\Gamma$. These linear contexts behave
% substructurally. As stated earlier, they are linear --- they do not obey
% weakening or contraction --- because a character is exhausted once read by a
% parsing procedure. Moreover, the characters in strings appear in the order in
% which we read them. We do not have the freedom to freely permute characters,
% therefore any type theory that is used to reason about formal grammars ought to
% omit the structural rule of exchange as well. This means that every variable
% within $\Delta$ must be used \emph{exactly once} and \emph{in order of
%   occurrence}. Thus, we can think of the linear contexts as an ordered list of
% limited resources. Once a resource is consumed, we cannot make reference to it
% again. Variables in a linear context then act like building
% blocks for constructing patterns over strings.

% We give the base types and type constructors for linear terms. As the
% interpretation of types as grammars in $\Grammar$ serves as our intended
% semantics, we simultaneously give the interpretation $\sem \cdot$ of the
% semantics as grammars.

% \pedro{The paragraphs above are a bit circular, what about the following alternative?}
% \steven{I am in favor of cutting what's above, but we haven't introduced the
%   syntax in the main body of the paper yet, even if the figures appeared
%   earlier. So what you write needs to take that into account}
% \pedro{In which case we should introduce it in the previous section. I suggest
%   adding this right after ``and the intuitionistic typing rules in fig 3''}
\subsection{Semantics}

We can now define our denotational semantics.
\begin{definition}[Grammar Semantics]
  We define the following interpretations by mutual recursion on the
  judgments of \theoryabbv:
  \begin{enumerate}
  \item For each non-linear context $\Gamma \isCtx$, we define a set $\sem \Gamma$.
  \item For each non-linear type $\Gamma \vdash X \isTy$, and element
    $\gamma \in \sem\Gamma$, we define a set $\sem X \gamma$.
  \item For each linear type $\Gamma A \isLinTy$ and element $\gamma
    \in \sem\Gamma$, we define a formal grammar $\sem{A}\gamma$. We
    similarly define a formal grammar $\sem\Delta\gamma$ for each
    linear context $\Gamma \Delta\isLinCtx$.
  \item For each non-linear term $\Gamma \vdash M : X$ and $\gamma \in \sem{\Gamma}$, we define an element $\sem{M}\gamma \in \sem{X}\gamma$.
  \item For each linear term $\Gamma; \Delta \vdash e : A$ and $\gamma \in \sem{\Gamma}$ we define a parse transformer from $\sem{\Delta}\gamma$ to $\sem{A}\gamma$.
  \end{enumerate}
  And we verify the following conditions:
  \begin{enumerate}
  \item If $\Gamma \vdash X \isSmall$, then $\sem X \gamma$ is a small set.
  \item If $\Gamma \vdash X \cong X'$ then for every $\gamma$, $\sem{X}\gamma = \sem{X'}\gamma$
  \item If $\Gamma \vdash A \cong A'$ then for every $\gamma$, $\sem{A}\gamma = \sem{A'}\gamma$
  \item If $\Gamma \vdash M \cong M' : X$ then for every $\gamma$, $\sem{M}\gamma = \sem{M'}\gamma$
  \item If $\Gamma;\Delta \vdash e \cong e' : A$ then for every $\gamma$, $\sem{e}\gamma = \sem{e'}\gamma$
  \end{enumerate}
\end{definition}

For reasons of brevity, we present only a selection of cases in
\Cref{fig:semantics}. The full details are in the appendix.\max{TODO:
  put the full details in the appendix}.
\max{TODO: how to describe $\mu$ succinctly?}\max{make sure $\lto$ and $\tol$ are right}
\begin{figure}
  \begin{align*}
    \sem{c}\gamma\,w &= \{ c | w = c\}\\
    \sem{I}\gamma\,w &= \{ () \pipe w = \varepsilon \}\\
    \sem{A \otimes B}\gamma\,w &= \{ (w_1,w_2,a,b) \pipe w_1w_2 = w \wedge a \in \sem{A}\gamma\,w_1 \wedge b \in \sem{B}\gamma\,w_2 \}\\
    \sem{\bigoplus_{x:X} A}\gamma\,w &= \{ (x, a) \pipe x \in \sem{X}\gamma \wedge a \in \sem{A}(\gamma,x)w \}\\
    \sem{\bigwith_{x:X} A}\gamma\,w &= \prod_{x\in \sem{X}\gamma} \sem{A}(\gamma, x)\,w\\
    \sem{A \lto B}\gamma\,w &= \prod_{w'} \sem{A}\gamma\,w' \to \sem{B}\gamma\,ww'\\
    \sem{B \tol A}\gamma\,w &= \prod_{w'} \sem{A}\gamma\,w' \to \sem{B}\gamma\,w'w\\
    \sem{\uparrow A}\gamma &= \sem{A}\gamma\,\varepsilon\\
    \sem{\SPF\,X}\gamma &= \{ F : \textrm{Functor}(\textrm{Grammar}^{\sem{X}\gamma}, \textrm{Grammar}) \pipe F \textrm{ preserves $\omega$-colimits.} \}\\
    \sem{\el(F)}\gamma\,G &= \sem{A}\gamma\,G\\
    \sem{\mu A}\gamma &= \mu (\sem{A}\gamma) \tag{\Cref{thm:adamek}}
  \end{align*}
  \caption{Grammar Semantics}
\end{figure}

\max{TODO: prose}

\max{TODO: spelling}
\begin{theorem}[Adamek]
  If a functor $F : \mathcal C \to \mathcal C$ preserves
  $\omega$-colimits and $\mathcal C$ has an initial object and
  $\omega$-colimits, then $F$ has an initial algebra given by the
  colimit of the chain
  \[ 0 \to F(0) \to F(F(0)) \to \cdots \]
\end{theorem}


%% That is, strings match $A \lto B$ if when prepended with a parse of $A$ they
%% complete to parses of $B$. In this manner, the linear function types generalize
%% Brzozowksi's notion of derivative
%% \cite{brzozowskiDerivativesRegularExpressions1964}.
%% Brzozowski initially only gave an accounting of this operation for
%% generalized regular expressions, but later work from Might et al.\ demonstrates that the same
%% construction can be generalized to context free grammars
%% \cite{mightParsingDerivativesFunctional2011}. Here, via the linear function
%% types, the same notion of derivative extends to the grammars of \theoryabbv.

%% Note, to ensure that the linear function types do indeed generalize Brzozowski
%% derivatives, we must include the equalities in \cref{fig:brzozowskideriv} as axioms\max{we must? why must we}.
%% \pedro{We should further explain why these equations are required, or at the very
%% least give some intuition}

%% Of course, all the above statements for the left function type also have
%% corresponding analogues for the
%% right-handed counterpart.

%% \begin{figure}
%% \begin{align*}
%%   c\lto c &\equiv I\\
%%   c\lto d &\equiv \bot\\
%%   c\lto I &\equiv \bot\\
%%   c\lto 0 &\equiv \bot\\
%%   c\lto (A \otimes B) & \equiv (c\lto A) \otimes B + (A \amp I) \otimes (c\lto B)\\
%%   c\lto A^* &\equiv (A \amp I)^{*} \otimes (c \lto A) \otimes A^* \\
%%   c\lto (A + B) &\equiv (c\lto A) + (c\lto B)
%% \end{align*}
%% \caption{Equality for Brzozowski Derivatives}
%% \label{fig:brzozowskideriv}
%% \end{figure}

%% \paragraph{LNL Dependent Types}
%% \steven{talk about how binary sum and binary with are definable over $\Bool$}
%% Given a non-linear type $X$, we may form both the dependent product type
%% $\LinPiTy x X A$ and the dependent pair type $\LinSigTy x X A$ as linear types
%% where $x$ is free in $A$.

%% $\sem{\LinPiTy x X A} \gamma w = \Pi(x:\sem{X}\gamma) \sem{A}(\gamma,x) w$

%% $\sem{\LinSigTy x X A} \gamma w = \Sigma(x:\sem{X}\gamma) \sem{A}(\gamma,x) w$

%% The grammar semantics of the linear product type is
%% indeed a dependent function out of the
%% semantics of $X$. Likewise, the grammar semantics of the linear dependent pair type is a dependent
%% pair in $\Set$.\pedro{The connection to grammar semantics needs to be better explained. For instance,
%% we should explain that disjunction corresponds to ``or'' of grammars}

%% Note that even though we do not take the binary additive conjunction and
%% disjunction as primitive, we may define them via these LNL dependent types. In
%% particular, via a dependence on $\Bool$.
%% %
%% \[
%%   A \amp B := \LinPiTy b \Bool {C(b)}
%% \]
%% \[
%%   A \oplus B := \LinSigTy b \Bool {C(b)},
%% \]
%% where $C(\true) = A, C(\false) = B$.

%% \paragraph{Universal Type}\max{this is absolutely not the right terminology. Universal type means multiple other completely different things}
%% \steven{Definable as nullary product}
%% The universal type $\top$ may be formed in any context.
%% %
%% \[ \sem{\top} \gamma w = \{ \ast \}\]
%% %
%% Its grammar semantics in set outputs the unit type in $\Set$ for all input strings in all contexts.

%% \paragraph{Empty Type}
%% \steven{Definable as nullary sum}
%% The empty type $\bot$ has no inhabitants. The elimination for the empty type
%% witnesses the principle of explosion --- i.e.\ from a term of type $\bot$ we may
%% introduce a term $\mathsf{absurd}_{A} : A$ for any type $A$.
%% %
%% \[ \sem {\bot} \gamma w = \emptyset \]
%% %
%% \paragraph{The $G$ Modality}
%% \steven{Rename this}
%% \steven{Compare to persistence modalities from separation logic}
%% The $G$ modality provides the embedding of linear types in the empty context
%% into non-linear types. The introduction and elimination forms for $G A$ obey the
%% same rules as given in \cite{bentonMixedLinearNonlinear1995} and \cite{krishnaswami_integrating_2015}.

%% The left adjoint $F$ from non-linear types back to linear types may be defined
%% as,
%% %
%% \[
%%   F X := \LinSigTy a A I
%% \]
%% %
%% \steven{Elaborate on $G$ more and rename $G A$to $\ltonl A$ in this section}

%% \pedro{Agreed :) maybe mention connections to persistent propositions from the
%% separation logic literature}

%% The semantics of $G A$ are exactly the semantics of $A$ at the empty word $\varepsilon$.
%% %
%% \[ \sem{G A} \gamma = \sem{A} \gamma \varepsilon \]
%% %
%% %
%% \paragraph{Recursive Types}
%% \steven{reflavor this section as ``indexed inductive'' instead of ``recursive'',
%% and entirely rework with the new inductive definition}
%% \steven{The non-indexed inductive are definable as being trivially indexed over unit}
%% Recursive linear types may be defined via the least-fixed point operator $\mu$.
% The grammar semantics of which are the fixed-point of sets induced by the
% grammar semantics of $A$.
% %
% \[ \sem{\mu x. A} \gamma = \mu (x:\Gr_i). \sem{A}(\gamma,x) \]
% %
% \pedro{We should justify, even if briefly, the existence of this
%   fixed-point}

% Observe that we need not take the Kleene star as a primitive
% grammar constructor, as it is definable as a fixed point.
% The Kleene star of a grammar $g$ is given as,

% \[
%   g^{*} := \mu X . I \oplus (g \otimes X)
% \]

% \begin{figure}[h!]
% \begin{mathpar}
%   \inferrule
%   {\Gamma ; \Delta \vdash p : I}
%   {\Gamma ; \Delta \vdash \mathsf{nil} : g^{*}}

%   \and

%   \inferrule
%   {\Gamma ; \Delta \vdash p : g \\ \Gamma ; \Delta' \vdash q
%   : g^{*}}
%   {\Gamma ; \Delta \vdash \mathsf{cons}(p , q) : g^{*}}

%   \\

%   \inferrule
%   {
%     \Gamma ; \Delta \vdash p : g^{*} \\
%     \Gamma ; \cdot \vdash p_{\varepsilon} : h \\
%     \Gamma ; x : g , y : h \vdash p_{\ast} : h
%   }
%   {\Gamma ; \Delta \vdash \mathsf{foldr}(p_{\varepsilon} , p_{\ast}) : g^{*}}
% \end{mathpar}
% \caption{Kleene Star Rules}
% \label{fig:star}
% \end{figure}

% Likewise, $g^{*}$ has admissible introduction and
% elimination rules, shown in \cref{fig:star}. Note that this
% definition of $g^{*}$ and these
% rules arbitrarily assigns a handedness to the Kleene star.
% We could have just as well took it to be a fixed point of
% $I \oplus (X \otimes g)$. In fact, the definitions are
% equivalent, as the existence of the $\mathsf{foldl}$ term below
% shows that $g^{*}$ is also a fixed point of
% $I \oplus (X \otimes g)$.
% This is a marked point of difference from Kleene
% algebra with recursion, where the fixed points for the left and right variants
% of Kleene star need not agree \cite{leiss}.

% \begin{equation}
%   \label{eq:foldl}
%   \inferrule
%   {
%     \Gamma ; \Delta \vdash p : g^{*} \\
%     \Gamma ; \cdot \vdash p_{\varepsilon} : h \\
%% %     \Gamma ; y : h , x : h \vdash p_{\ast} : h
%% %   }
%% %   {\Gamma ; \Delta \vdash \mathsf{foldl}(p_{\varepsilon} , p_{\ast}) : g^{*}}
%% % \end{equation}

%% % In fact, the $\mathsf{foldl}$ term is defined using
%% % $\mathsf{foldr}$ --- much in the same way one
%% % may define a left fold over lists in terms of a right fold
%% % in a functional programming language.
%% % The underlying trick is to fold over a list of linear functions
%% % instead of the original string. We curry each character $c$
%% % of the string into a function that concatenates $c$, and
%% % right fold over this list of linear functions. Because function
%% % application is left-associative, this results in a left
%% % fold over the original string.

%% % We only take fixed points of a single variable as a
%% % primitive operation in the type theory, but we may apply
%% % Beki\`c's theorem \cite{Bekić1984} to define an admissible
%% % notion of multivariate fixed point. This is particularly
%% % useful for defining grammars that encode the states of an
%% % automaton, as we will see in \cref{sec:automata}. In \cref{fig:multifix} we provide the
%% % introduction and elimination principles for such a fixed
%% % point, where $\sigma$ is the substitution that unrolls the
%% % mutually recursive definitions one level. That is,

%% % \begin{align*}
%% %   \sigma = \{ & \mu(X_{1} = A_{1} \dots, X_{n} = A_{n}).X_{1} / X_{1} , \dots, \mu(X_{1} = A_{1}, \dots, X_{n} = A_{n}).X_{n} / X_{n} \}
%% % \end{align*}

%% % \begin{figure}
%% % \begin{mathpar}
%% %   \inferrule
%% %   {\Gamma ; \Delta \vdash e : \simulsubst {A_{k}} {\sigma}}
%% %   {\Gamma ; \Delta \vdash \mathsf{cons}~e : \mu(X_{1} = A_{1}, \dots, X_{n} = A_{n}).X_{k}}

%% %   \and

%% %   \inferrule
%% %   {\Gamma ; \Delta \vdash e : \mu(X_{1} = A_{1}, \dots, X_{n} = A_{n}).X_{k} \\
%% %              \Gamma ; x_{j} : \simulsubst {A_{j}}{\sigma} \vdash e_{j} : B_{j}\quad \forall j
%% %   }
%% %   {\Gamma; \Delta \vdash \mathsf{mfold}(x_{1}.e_{1}, \dots, x_{n}.e_{n})(e) : B_{k}}
%% % \end{mathpar}
%% % \caption{Multi-fixed Points}
%% % \label{fig:multifix}
%% % \end{figure}

%% %% \section{Categorical Semantics of \theoryabbv}
%% %% \label{sec:categorify}

%% As described in previous sections, \theoryabbv has a ``standard''
%% interpretation where non-linear types denote sets, linear types denote
%% formal grammars, non-linear terms denote functions and linear terms
%% denote parse transformers.
%% %
%% Further, as is typical for type theory, every type constructor in the
%% calculus is interpreted by using a \emph{universal construction} in
%% either the category of sets or the category of grammars\footnote{the
%% only exception is that the axioms for Brzozowski derivatives we
%% considered do not follow solely from universal properties. In this
%% section we focus on models that do not necessarily satisfy these
%% axioms.}.
%% %
%% This leads to an immediate opportunity for generalization: the type
%% theory has in addition to the standard grammar-theoretic
%% interpretation, an interpretation in any categorical structure that
%% exhibits the same universal constructions.
%% %
%% On the one hand this gives us a structured way to formalize the
%% standard semantics precisely, but additionally it enables us to
%% consider \emph{non-standard} models in the next section that point to
%% further applications as well as meta-theoretic results.

%% In developing this categorical semantics, we will start with the
%% closest analogue from formal language theory: Kleene algebra.  Kleene
%% algebras are an important tool in the theory of regular languages
%% serving as a bridge between algebraic reasoning and equivalence of
%% regular expressions. More broadly, through various extensions, they
%% serve as a theoretical substrate to studying different kinds of formal
%% languages.
%% %
%% We can then see that the ``regular fragment'' of our type theory
%% (i.e., just characters, $\otimes$, $\oplus$, $0,1$ and a Kleene star)
%% has models in what we call a \emph{Kleene category} a categorification
%% of Kleene algebra from posets to categories.
%% %
%% We then further develop this into our final notion of model, which we
%% call a \emph{Chomsky category} as it can model not just regular
%% languages, but the full Chomsky hierarchy.

%% \subsection{Comparison to Kleene Algebra}
%% \subsection{Kleene Algebra and Kleene Categories}
%% A Kleene algebra is a tuple $(A, +, \cdot, (-)^*, 1, 0)$, where $A$ is
%% a set, $+$ and $\cdot$ are binary operations over $A$, $(-)^*$ is a
%% function over $A$, and $1$ and $0$ are constants. These structures
%% satisfy the axioms depicted in Figure~\ref{fig:axioms}.

%% \begin{figure}
%%   \begin{align*}
%%     x + (y + z) &= (x + y) + z & x + y &= y + x\\
%%     x + 0 &= x & x + x &= x\\
%%     x(yz) &= (xy)z & x1 &= 1x = x\\
%%     x(y + z) &= xy + xz & (x + y)z &= xz + yz\\
%%     x0 &= 0x = x & & \\
%%     1 + aa^* &\leq a^* & 1 + a^*a &\leq a^*\\
%%      b + ax \leq x &\implies a^*b \leq x &  b + xa \leq x &\implies ba^* \leq x
%%   \end{align*}
%%   \caption{Kleene algebra axioms}
%%   \label{fig:axioms}
%% \end{figure}

%% The addition operation can be used to define the partial order
%% structure $a \leq b$ if, and only if, $a + b = b$. In the theory of
%% formal languages, this order structure can be used to model language
%% containment. In this section, we categorify the concept of Kleene
%% algebra and build on top of it in order to define an abstract theory
%% of parsing. We start by defining \emph{Kleene categories}.

%% \begin{definition}
%%   A Kleene category is a distributive monoidal category $\cat{K}$
%%   such that for every objects $A$ and $B$, the endofunctors $F_{A, B}
%%   = B + A \otimes X$ and $G_{A, B} = B + X \otimes A$ have initial
%%   algebras (denoted $\mu X.\, F_{A, B}(X)$) such that $B \otimes (\mu
%%   X.\, F_{A, 1}) \cong \mu X.\, F_{A, B}(X)$ and the analogous isomorphism
%%   for $G_{A,B}$ also holds.
%% \end{definition}

%% As a sanity check, note that Kleene algebras are indeed examples of
%% Kleene categories.

%% \begin{example}
%%   Every Kleene algebra, seen a posetal category, is a Kleene category.
%%   The product $\cdot$ is a monoidal product and the addition is a
%%   least-upper bound, which corresponds to a coproduct. Lastly, the
%%   axioms of the Kleene star have a direct correspondence to the
%%   coherence conditions postulated by the initial algebras of Kleene
%%   categories.
%% \end{example}

%% This example provides a neat categorical justification to how
%% restrictive Kleene algebras are in terms of reasoning about
%% languages. By only having at most one morphism between objects, there
%% is not a lot of information they can convey. In this case, the only
%% information you get is language containment. As demonstrated by
%% $\mathbf{Gr}$, the extra degrees of freedom granted by having more
%% morphisms gives you more algebraic structure for reasoning about
%% languages.

%% For the next example, we see an unexpected connection with the theory
%% of substructural logics.

%% \begin{example}
%%   The opposite category of every Kleene category is a model of a variant of
%%   conjunctive ordered logic, where the Kleene star plays the role of the ``of
%%   course'' modality from substructural logics which allows hypotheses to
%%   be discarded or duplicated.
%% \end{example}

%% As we have seen, the proposed axioms are a direct translation of the
%% Kleene algebra axioms to a categorical setting. Its most unusual aspect is the
%% axiomatization of the Kleene star as a family of initial algebras
%% satisfying certain isomorphisms. If the Kleene category $\cat{K}$ has
%% more structure, then these isomorphisms hold ``for free''.

%% \begin{theorem}
%%   \label{th:kleeneclosed}
%%   Let $\cat{K}$ be a Kleene category such that it is also monoidal
%%   closed.  Then, the initial algebras isomorphisms hold automatically.
%% \end{theorem}
%% \begin{proof}
%%   We prove this by the unicity (up-to isomorphism) of initial
%%   algebras. Let $[hd, tl]: I + (\mu X.\, F_{A, I}(X)) \otimes A \to
%%   (\mu X.\, F_{A, I}(X))$ be the initial algebra structure of $(\mu
%%   X.\, F_{A, I}(X))$ and consider the map $[hd, tl] : B + B \otimes
%%   (\mu X.\, F_{A, I}(X)) \otimes A \to B\otimes (\mu X.\, F_{A,
%%     I}(X))$.

%%   Now, let $[f,g] : B + A \otimes Y \to Y$ be an $F_{A,B}$-algebra and
%%   we want to show that there is a unique algebra morphism $h : \mu X.\, F_{A,I} \to B \lto Y$. We can show existence and
%%   uniqueness by showing that the diagram on top commutes if, and
%%   only if, the diagram on the bottom commutes:

%% % https://q.uiver.app/#q=WzAsOCxbMCwwLCJCICsgQiBcXG90aW1lcyAoXFxtdSBYLlxcLCAxICsgWCBcXG90aW1lcyBBKSJdLFsyLDAsIkIgKyBZIFxcb3RpbWVzIEEiXSxbMiwyLCJZIl0sWzAsMiwiQiBcXG90aW1lcyAoXFxtdSBYLlxcLCAxICsgWCBcXG90aW1lcyBBKSJdLFswLDMsIjEgKyAoXFxtdSBYLlxcLCAxICsgWCBcXG90aW1lcyBBKSJdLFswLDUsIlxcbXUgWC5cXCwgMSArIFggXFxvdGltZXMgQSJdLFsyLDMsIjEgKyAoIEIgXFxsdG8gWSkgXFxvdGltZXMgQSJdLFsyLDUsIkIgXFxsdG8gWSJdLFswLDEsImlkICsgKGlkIFxcb3RpbWVzIGg7IGV2KSBcXG90aW1lcyBpZF9BIl0sWzEsMiwiW2YsZ10iXSxbMywyLCJpZCBcXG90aW1lcyBoOyBldiIsMl0sWzAsMywiW2lkIFxcb3RpbWVzIGgsIGlkIFxcb3RpbWVzIHRsXSIsMl0sWzQsNiwiaWQgKyAoaCBcXG90aW1lcyBYKVxcb3RpbWVzIGlkIl0sWzUsNywiaCIsMl0sWzYsNywiW2YnLCBnJ10iXSxbNCw1LCJbaGQsIHRsXSIsMl1d
%% \[\begin{tikzcd}
%% 	{B + B \otimes (\mu X.\, I + X \otimes A)} && {B + Y \otimes A} \\
%% 	\\
%% 	{B \otimes (\mu X.\, I + X \otimes A)} && Y \\
%% 	{I + (\mu X.\, I + X \otimes A)} && {I + ( B \lto Y) \otimes A} \\
%% 	\\
%% 	{\mu X.\, I + X \otimes A} && {B \lto Y}
%% 	\arrow["{id + (id \otimes h; ev) \otimes id_A}", from=1-1, to=1-3]
%% 	\arrow["{[f,g]}", from=1-3, to=3-3]
%% 	\arrow["{id \otimes h; ev}"', from=3-1, to=3-3]
%% 	\arrow["{[id \otimes h, id \otimes tl]}"', from=1-1, to=3-1]
%% 	\arrow["{id + (h \otimes X)\otimes id}", from=4-1, to=4-3]
%% 	\arrow["h"', from=6-1, to=6-3]
%% 	\arrow["{[f', g']}", from=4-3, to=6-3]
%% 	\arrow["{[hd, tl]}"', from=4-1, to=6-1]
%% \end{tikzcd}\]
%%   This equivalence follows by using the adjunction structure given
%%   by the monoidal closed structure of $\cat{K}$. A completely analogous
%%   argument for $G_{A,B}$ also holds. Furthermore, by generalizing the
%%   construction of \Cref{sec:formaltype}, we can also show that from
%%   the monoidal closed assumption it follows that $\mu X.\, F_{A, I}(X) \cong \mu X.\, G_{A, I}(X)$
%% \end{proof}

%% Something surprising about this lemma is that it provides an alternative
%% perspective on the observation that if a Kleene algebra has an
%% residuation operation, also called action algebra \cite{kozen1994action},
%% then the Kleene star admits a simpler axiomatization.

%% Since we want Kleene categories to generalize our notion of formal
%% grammars as presheaves $\String \to \Set$, we prove that they do
%% indeed form a Kleene category. We start by presenting a well-known
%% construction from presheaf categories.

%% \begin{definition}
%%   Let $\cat{C}$ be a locally small monoidal category and $F$, $G$ be
%%   two functors $\cat{C} \to \Set$. Their Day convolution tensor
%%   product is defined as the following coend formula:
%%   \[
%%   (F \otimes_{Day} G)(x) = \int^{(y,z) \in \cat{C}\times\cat{C}}\cat{C}(y\otimes z, x) \times F(y) \times G(z)
%%   \]
%%   Dually, its internal hom is given by the following end formula:
%%   \[
%%   (F \lto_{Day} G)(x) = \int_{y} \Set(F(y), G(x \otimes y))
%%   \]
%% \end{definition}

%% \begin{lemma}[Day \cite{day1970construction}]
%%   Under the assumptions above, the presheaf category $\Set^{\cat{C}}$ is
%%   monoidal closed.
%% \end{lemma}

%% %% \begin{theorem}
%% %%   Let $\cat{K}$ be a Kleene category and $A$ a discrete category.
%% %%   The functor category $[A, \cat{K}]$.
%% %%   (HOW GENERAL SHOULD THIS THEOREM BE? BY ASSUMING ENOUGH STRUCTURE,
%% %%   E.G. K = Set, THIS THEOREM BECOMES SIMPLE TO PROVE)
%% %% \end{theorem}
%% \begin{theorem}
%%   If $\cat{C}$ is a locally small monoidal category, then
%%   $\Set^{\cat{C}}$ is a Kleene category.
%% \end{theorem}
%% \begin{proof}

%%   By the lemma above, $\Set^{\cat{C}}$ is monoidal closed, and since it
%%   is a presheaf category, it has coproducts. Furthermore, the tensor
%%   is a left adjoint, i.e. it preserves colimits and, therefore, it is
%%   a distributive category.

%%   As for the Kleene star, since presheaf categories admit small colimits,
%%   the initial algebra of the functors $F_{A,B}$ and $G_{A,B}$ can be
%%   defined as the filtered colimit of the diagrams:

%%   From Theorem~\ref{th:kleeneclosed} it follows that these initial
%%   algebras satisfy the required isomorphisms and this concludes the
%%   proof.
%% \end{proof}

%% \begin{corollary}
%%   For every alphabet $\Sigma$, the presheaf category $\Set^{\cat{\Sigma^*}}$
%%   is a Kleene category.
%% \end{corollary}
%% \begin{proof}
%%   Note that string concatenation and the empty string make the
%%   discrete category $\Sigma^*$ a strict monoidal category.
%% \end{proof}

%% Much like in the posetal case, the abstract structure of a Kleene
%% category is expressive enough to synthetically reason about regular
%% languages. A significant difference between them is that while Kleene
%% algebras can reason about language containment, Kleene categories can
%% reason about \emph{ambiguity}, \emph{strong equivalence} of grammars.

%% \subsection{Lambek Hyperdoctrines and Chomsky Hyperdoctrines}

%% Though Kleene categories are expressive enough to reason about
%% concepts that are outside of reach of Kleene algebras, their
%% simply-typed nature makes them not so expressive from a type theoretic
%% point of view. This is limiting because type theories are successful
%% syntactic frameworks for manipulating complicated categorical
%% structures while avoiding some issues common in category theory, such
%% as coherence issues.

%% With this in mind, we want to design a categorical semantics that
%% builds on top of Kleene categories with the goal of extending them
%% with dependent types and making them capable reasoning about languages
%% and their parsers. This leads us to the abstract notion of model we
%% are interested in capturing with \theoryabbv: a \emph{Chomsky
%% category}.

%% We do this in two stages: first we define a \emph{Lambek hyperdoctrine} to be a
%% notion of model for the judgmental structure of \theoryabbv: that is we have
%% notions of linear and non-linear type, contexts, substitution and terms, but do
%% not assume any particular type constructions exist. Then we will define when a
%% Lambek hyperdoctrine is a Chomsky hyperdoctrine, meaning it can interpret all
%% the type formers of \theoryabbv and therefore arbitrary strength grammars.

%% \begin{definition}[Lambek Hyperdoctrine]
%%   A \emph{Lambek hyperdoctrine} consists of
%%   \begin{enumerate}
%%   \item A category $\mathcal C$ with a terminal object.
%%   \item A category with families structure over $\mathcal C$.
%%   \item A contravariant functor $L : \mathcal C^{o} \to \textrm{MultiCat}$ from
%%     $\mathcal C$ to the category of Multicategories.
%%   \end{enumerate}
%% \end{definition}
%% Here the objects of $\mathcal C$ model the non-linear contexts, and morphisms
%% model non-linear substitutions. The category with families structure over
%% $\mathcal C$ models the dependent non-linear types. Finally the
%% ``hyperdoctrine'' of multicategories $L$ models the linear types and terms. The
%% fact that this is functorial in $\mathcal C$ corresponds to the fact that linear
%% types are all relative to a non-linear context and that linear variables and
%% composition commute with non-linear substitution.

%% We will have three main Lambek hyperdoctrines of interest in this paper: the
%% ``standard model'' of sets and grammars, the ``syntactic model'' given by our
%% type theory itself and lastly a gluing model introduced in
%% Section~\ref{sec:canonicity} to prove the canonicity theorem.
%% \begin{example}
%%   \begin{enumerate}
%%   \item Let $\mathcal C$ be the category of sets, equipped with its usual
%%     category with families structures of families. Define $L : \mathcal C^{o}
%%     \to \textrm{MultiCat}$ to map $L(X)$ to the representable multicategory
%%     $(\Set^{\Sigma^*})^X$ of $X$-indexed families of grammars where the monoidal
%%     structure is given pointwise by the Day convolution monoidal structure.
%%   \item We can build a ``syntactic'' Lambek hyperdoctrine $\Syn(\Sigma)$ from the syntax
%%     itself: the category $\mathcal C$ is given by non-linear contexts and
%%     substitutions, the category with families by the non-linear types and terms
%%     and the hyperdoctrine of multicategories by the linear types and terms.
%%   \end{enumerate}
%% \end{example}

%% Next each type linear and non-linear type constructor corresponds to extra data
%% on a Lambek hyperdoctrine.
%% \begin{definition}
%%   \label{def:chomsky-data}
%%   \begin{enumerate}
%%   \item Dependent type structure: standard (say extensional type theory, give a reference)
%%   \item Universes
%%   \item Inductive grammars?
%%   \item Non-inductive Grammar constructors: standard monoidal category stuff
%%   \end{enumerate}
%% \end{definition}
%% \begin{definition}[Chomsky Hyperdoctrine]
%%   A Chomsky Hyperdoctrine is a Lambek hyperdoctrine equipped with all of the data in Definition~\ref{def:chomsky-data}.
%% \end{definition}
%% \pedro{We should probably define some of the words in this definition}
%% \begin{definition}
%%   A Chomsky category is a locally Cartesian category with two hierarchies of
%%   universes $\{L_i\}_{i\in \nat}$ and $\{U_i\}_{i\in \nat}$ such that
%%   every $L_i$ and $U_i$ are $U_{i+1}$-small. Furthermore, we require
%%   $U_i$ to be closed under dependent products and sum,
%%   $L_i$ to be closed under the Kleene category connectives,
%%   dependent products, left and right closed structures, with
%%   a type constructor $G : L_i \to U_i$ and a linear dependent sum
%%   going the other direction.
%% \end{definition}

%% \steven{Max suggests augmenting the definition of a Chomsky category to
%%   something like two categories $L$ and $U$, and $U$ is Cartesian closed I don't
%%   fully recall the rest.


%%   My best guess is that you take $L$ a Kleene category with a hierarchy of universes
%%   two, and you further require that $L$ is $Psh(U)$-enriched, except he
%%   suggested a further adjective on these presheaves. Perhaps representability?
%% }

%% \begin{theorem}
%%   The presheaf category $\Grammar = \Set^{\cat{\Sigma^*}}$ is a Chomsky category.
%% \end{theorem}


%% Further, the syntactic category of \theoryname is manifestly a Chomsky category.

%% \pedro{This is likely true, but if we explicitly say so, this warrants a proof. I think that
%% if we don't say anything about the syntactic category, reviewers won't mind.}

%% \steven{I agree that we don't want to say anything that opens unnecessary
%%   questions for proofs we haven't written. However, it seems hard to make the
%%   case that we have the right categorical model of the syntax if this statement
%%   isn't true. By restating our definition of Chomsky category, this should be
%%   obvious or a quick proof}

%% \pedro{I agree, then we should add the quick proof :) }

%% \subsection{Concrete Models of \theoryabbv}
%% \label{sec:othermodels}

%% \steven{Because we can define a version of semantic actions internally, we
%%   shouldn't put it as a separate model}

%% One of the powers of type theories is that it can be profitable to interpret them
%% in various models. In this section, by using our just-defined Chomsky categories,
%% we show how other useful concepts from formal language theory can also be organized
%% as models of \theoryabbv. We illustrate this point by providing two examples that
%% are closely related to the theory of formal languages: language equivalence and
%% semantic actions. Furthermore, in order to justify how $\mathbf{Gr}$ relates to
%% more traditional notions of parsing, we define a glued model that proves a
%% canonicity property of grammar terms.


%% \subsection{Language Semantics}
%% Every grammar induces a language semantics. Also languages can be taken as a
%% propositionally truncated view of the syntax. Logical equivalence should induce
%% weak equivalence, and thus even give a syntactic way to reason about language equivalence.
%% \steven{TODO language semantics}

%% % \subsection{Semantic Actions}
%% % \steven{Tentatively planning to cut this subsection for an internal representation
%% %   of semantic actions}
%% % Returning to the problem of parsing, the output of a parse usually is not the
%% % literal parse tree. Rather, the output is the result of some \emph{semantic
%% %   actions} ran on the parse tree, which usually serve to remove some syntactic
%% % details that are unnecessary for later processing.

%% % Given a grammar $G : \String \to \Set$, we define a semantic action to be a set
%% % $X$ with a function $f$ that produces a semantic element from any parse of $G$.

%% % \[
%% %   f : \PiTy w \String {G w \to X}
%% % \]

%% % Further, semantic actions can be arranged into a structured category.
%% % Define $\SemAct$ as the comma
%% % category $\Grammar / \Delta$, where $\Delta : \Set \to \Grammar$ defines a
%% % discrete presheaf. That is, for a set $X$, $\Delta (X)(w) = X$ for all
%% % $w \in \String$. As $\SemAct$ is defined as a comma category, it has a forgetful
%% % functor into $\Grammar$. That is, $\SemAct$ serves as a notion of formal
%% % grammar. Moreover, $\SemAct$ is a model of \theoryabbv.

%% % \steven{It being a notion of formal grammar is distinct from being a model. This
%% % probably warrants a proof}

%% % \steven{TODO semantic actions}

%% % \pedro{This is a very nice opportunity of showing off the supremacy of denotational
%% %   reasoning ;) We should probably prove the gluing lemma in the previous section
%% %   and apply it here and in the canonicity section. The actual proof might have to be
%% %   moved to the appendix, though}

%% \subsection{Parse Canonicity}
%% Canonicity is an important metatheoretic theorem in the type theory
%% literature.  It provides insight on the normal forms of terms and,
%% therefore, on its computational aspects. Frequently, proving
%% canonicity for boolean types, i.e. every closed term of type bool
%% reduces to either true or false, is enough to justify that the type
%% theory being studied is well-behaved. In our case, however, since we
%% want to connect \theoryabbv to parsers, we must provide a more
%% detailed account of canonicity. In particular, we give a nonstandard semantics
%% of \theoryabbv that carries a proof of canonicity along with it.

%% If $\cdot \vdash A$ is a closed linear type then there are
%% two obvious notions of what constitutes a ``parse'' of a string w
%% according to the grammar $A$:
%% \begin{enumerate}
%% \item On the one hand we have the set-theoretic semantics just
%%   defined, $\llbracket A \rrbracket \cdot w$
%% \item On the other hand, we can view the string $w = c_1c_2\cdots$ as
%%   a linear context $\lceil w \rceil = x_1:c_1,x_2:c_2,\ldots$ and
%%   define a parse to be a $\beta\eta$-equivalence class of linear terms $\cdot;
%%   \lceil w \rceil \vdash e : G$.
%% \end{enumerate}
%% It is not difficult to see that at least for the ``purely positive''
%% formulae (those featuring only the positive connectives
%% $0,+,I,\otimes,\mu, \overline\Sigma,c$) that
%% every element $t \in \llbracket A \rrbracket w$ is a kind of tree and
%% that the nodes of the tree correspond precisely to the introduction
%% forms of the type. However it is far less obvious that \emph{every}
%% linear term $\lceil w \rceil \vdash p : \phi$ is equal to some
%% sequence of introduction forms since proofs can include elimination
%% forms as well. To show that this is indeed the case we give a
%% \emph{canonicity} result for the calculus: that the parses for .

%% \begin{definition}
%%   A non-linear type $X$ is purely positive if it is built up using
%%   only finite sums, finite products and least fixed points.

%%   A linear type is purely positive if it is built up using only finite
%%   sums, tensor products, generators $c$, least fixed points and linear
%%   sigma types over purely positive non-linear types.
%% \end{definition}

%% \begin{definition}
%%   %% Let $X$ be a closed non-linear type. The closed elements $\textrm{Cl}(X)$ of $X$ are the definitional equivalence classes of terms $\cdot \vdash e : X$.

%%   Let $A$ be a closed linear type. The nerve $N(A)$ is a presheaf on
%%   strings that takes a string $w$ to the definitional equivalence
%%   classes of terms $\cdot; \lceil w\rceil \vdash e: N(A)$.
%% \end{definition}

%% \begin{theorem}[Canonicity]
%%   Let $A$ be a closed, purely positive linear type. Then there is an
%%   isomorphism between $\llbracket A\rrbracket$ and $N(A)$.
%% \end{theorem}
%% \begin{proof}
%%   We outline the proof here, more details are in the appendix. The
%%   proof proceeds first by a standard logical families construction
%%   that combines canonicity arguments for dependent type theory
%%   TODO cite coquand
%%   % \cite{coquand,etc}
%%   with logical relations constructions for linear
%%   types
%%   TODO cite hylandschalk
%%   % \cite{hylandschalk}
%%   . It is easy to see by induction that the
%%   logical family for $A$, $\hat A$ is isomorphic to $\llbracket A
%%   \rrbracket$ and the fundamental lemma proves that the projection
%%   morphism $p : \hat A \to N(A)$ has a section, the canonicalization
%%   procedure. Then we establish again by induction that
%%   canonicalization is also a retraction by observing that introduction
%%   forms are taken to constructors.
%% \end{proof}


%% \begin{enumerate}
%% \item Every term $\lceil w \rceil \vdash p : G + H$ is equal to $\sigma_1q$ or $\sigma_2 r$ (but not both)
%% \item There are no terms $\lceil w \rceil \vdash p : 0$
%% \item If there is a term $\lceil w \rceil \vdash p : c$ then $w = c$ and $p = x$.
%% \item Every term $\lceil w \rceil \vdash p : G \otimes H$ is equal to $(q,r)$ for some $q,r$
%% \item Every term $\lceil w \rceil \vdash p : \varepsilon$ is equal to $()$
%% \item Every term $\lceil w \rceil \vdash p : c$ is equal to $x:c$
%% \item Every term $\lceil w \rceil \vdash p : \mu X. G$ is equal to $\textrm{roll}(q)$ where $q : G(\mu X.G/X)$
%% \item Every term $\lceil w \rceil \vdash p : (x:A) \times G$ is equal
%%   to $(M,q)$ where $\cdot \vdash M : A$
%% \end{enumerate}

%% To prove this result we will use a logical families model. We give a
%% brief overview of this model concretely:
%% \begin{enumerate}
%% \item A context $\Gamma$ denotes a family of sets indexed by closing substitutions $\hat\Gamma : (\cdot \vdash \Gamma) \Rightarrow \Set_i$
%% \item A type $\Gamma \vdash X : U_i$ denotes a family of sets $\hat X : \Pi(\gamma:\cdot \vdash \Gamma) \hat\Gamma \Rightarrow (\cdot \vdash \simulsubst X \gamma) \Rightarrow \Set_i$
%% \item A term $\Gamma \vdash e : X$ denotes a section $\hat e : \Pi(\gamma)\Pi(\hat\gamma)\hat X \gamma \hat\gamma (\simulsubst e \gamma)$
%% \item A linear type $\Gamma \vdash A : L_i$ denotes a family of grammars $\hat A : \Pi(\gamma:\cdot\vdash\Gamma)\,\hat\Gamma \Rightarrow \Pi(w:\Sigma^*) (\cdot;\lceil w\rceil \vdash A[\gamma])\Rightarrow \Set_i$, and the denotation of a linear context $\Delta$ is similar.
%% \item A linear term $\Gamma;\Delta \vdash e : A$ denotes a function \[\hat e : \Pi(\gamma)\Pi(\hat\gamma)\Pi(w)\Pi(\delta : \lceil w \rceil \vdash \simulsubst \Delta \gamma) \hat\Delta \gamma \hat\gamma \delta \Rightarrow \hat A \gamma \hat\gamma w {(\simulsubst {\simulsubst e \gamma} \delta)}\]
%% \end{enumerate}
%% And some of the constructions:
%% \begin{enumerate}
%% \item $\widehat {(G A)} \gamma \hat\gamma e = \hat A \gamma \hat\gamma \varepsilon (G^{-1}e)$
%% \item $\widehat {(A \otimes B)} \gamma \hat\gamma w e = \Sigma(w_Aw_B = w)\Sigma(e_A)\Sigma(e_B) (e_A,e_B) = e \wedge \hat A \gamma \hat \gamma w_A e_A \times \hat B \gamma \hat \gamma w_B e_B$
%% \item $\widehat {(A \lto B)} \gamma \hat\gamma w e = \Pi(w_A)\Pi(e_A) \hat A \gamma\hat\gamma w_A e_A \Rightarrow \hat B \gamma\hat\gamma (ww_A) (\applto {e_A} e)$
%% \end{enumerate}

%% First, the category with families will be
%% the category of logical families over set contexts/types
%% $\Delta$/$A$. Then the propositional portion will be defined by
%% mapping a logical family $\hat \Gamma \to \Gamma$

%% First, let $L$ be the category of BI formulae and proofs (quotiented
%% by $\beta\eta$ equality). Define a functor $N : L \to \Set^{\Sigma^*}$ by
%% \[ N(\phi)(w) = L(w,\phi) \]

%% Then define the gluing category $\mathcal G$ as the comma category
%% $\Set^{\Sigma^*}/N$. That is, an object of this category is a pair of
%% a formula $\phi \in L$ and an object $S \in \mathcal
%% P(\Sigma^*)/N(\phi)$. We can then use the equivalence $\mathcal
%% P(\Sigma^*)/N(\phi) \cong \mathcal P(\int N(\phi))$ to get a simple
%% description of such an $S$: it is simply a family of sets indexed by
%% proofs $L(w,\phi)$:
%% \[ \prod_{w\in\Sigma^*} L(w,\phi) \to \Set \]
%% This category clearly comes with a projection functor $\pi : \mathcal
%% G \to \mathcal L$ and then our goal is to define a section by using
%% the universal property of $\mathcal L$.

%% To this end we define
%% \begin{enumerate}
%% \item $(\phi, S) \otimes (\psi, T) = (\phi \otimes \psi, S\otimes T)$ where
%%   \[ (S \otimes T)(w, p) = (w_1w_2 = w) \times (q_1,q_2 = p) \times S\,w_1\,q_1 \times T\,w_2\,q_2\]
%% \item $(\phi, S) \multimap (\psi, T) = (\phi \multimap \psi, S \multimap T)$ where
%%   \[ (S \multimap T)(w,p) = w' \to q \to S\,w'\,q \to T (ww') (p\,q) \]
%% \item $\mu X. ??$ ??
%% \end{enumerate}

%% \pedro{We should conclude this section by explaining the relevance of the canonicity
%% theorem. Could also be done before stating the theorem.}

\section{Discussion and Future Work}
\label{sec:discussion}
\paragraph{Lambek Calculus and Categorial Grammar}
In 1958 \cite{lambek1958mathematics}, Joachim Lambek introduced his syntactic
calculus as a logical system for linguistic derivations. In fact, one may view
the subtheory of our calculus generated over $\otimes, \lto, \tol$ to be a proof
relevant presentation of Lambek's original system.

\steven{Should discussion of this connection to Lambek calculus be spread
  throughout the paper?}

Lambek later explicitly connects his syntactic calculus to the structure of a
non-commutative biclosed monoidal category \cite{lambek1988categorial}. There is
a rich line of subsequent research after these revelations by Lambek, including
the study of pregroup grammars
\cite{coeckeMathematicalFoundationsCompositional2010}, the presentation of
abstract categorial grammars \cite{degrooteAbstractCategorialParsing2015}, and
even Lambek calculus with dependent types \cite{luoLambekCalculusDependent}. As
far as the authors can tell, these works seem to be primarily interested in natural
language semantics in the style of Montague grammars, rather than parsing formal
languages. That is, they seem to try to give a functorial semantic
interpretation of natural language.

There seem to be some implementations from the abstract categorial grammar
toolkit, but again they seem natural language focused and not up to the task of
verified formal language parsing.

Also mention here that \cite{luoLambekCalculusDependent} seems similar at first
glance, but isn't as similar as it might look because the judgments and
dependent types seem to
take on a different interpretation of ours. Say something here like, ``as we
understand it, it differs in x,y,z ways''

\steven{TODO Edit the above}
\steven{Even though the above sounds bad, are there Lambek/categorial things that are missed?}

%% 1. More practical: integrating it into a larger verified development
%%    - verified imperative implementation: Lambek logic ala separation
%%      logic?
%% 2. Semantic actions
%% 3. Type systems as tree grammars

\paragraph{Typed Approaches to Regular Expressions}
The usage of a simple type system to reason about regular expressions was
introduced by Frisch and Cardelli in 2004 \cite{frischCardelli}, and later expanded
by Hengelin and Nielsen in 2011 \cite{henglein_regular_2011} to handle proofs of regular expression
containment. In 1992, while investigating the addition of arbitrary recursion to Kleene
algebra \cite{leiss},
Lei{\ss} used a least fixed-point operator instead of the Kleene star. Although
he did not explicitly make use of a type system, or have access to the view of Frisch and
Cardelli, Lei{\ss} had added the full power of inductive types to a type system
of regular expressions.

Pratt's work on the residuation of action algebras in \cite{prattActionLogicPure1991} closely mirrors
our observation that Brzozowski derivatives mirror linear function types.
Moreover, he shows that his residuals form a Galois connection with sequential
composition. This is precisely a posetal --- or \emph{thin} --- version of the
adjunction between
linear function types and tensor product in \theoryabbv.

Our primary observation in this paper can be summarized as saying that all of
these above works can embed naturally into a rich, unifying theory.

\steven{I'd like to synthesize the discussion of these works with the categorial
grammar works, but I am not sure how. It seems as though there are these
parallel tracks of research that have missed each other for decades}


% In \cite{firsovCertifiedNormalizationContextFree2015}, Firsov and Uustalu
% constructed a verified normalization procedure to turn a
% context-free grammar into an equivalent grammar in Chomsky normal form.
% Combining their work with a CYK parser, they have built a verified context-free
% grammar parser, \emph{up to weak equivalence}. To upgrade from a weak
% equivalence to a strong equivalence, they suggest updating their grammar
% formalism to treat parse trees as first-class proofs of language membership.
% That is, without explicitly stating so, they suggest that they next natural step
% is to reason about grammars in a system where parse trees are terms, as we have
% given in this paper.

% \steven{Make the above paragraph shorter. Also idk if its worth including, but I
% like that they practically ask for a type theory}

\steven{Mention Neel's typed algebraic approach paper here. Say it is similar at
  first glance, as its parsing with types, but that he actually is representing
  something quite different. Or is it so different that its not worth mentioning?}

\steven{TODO add Lambek calculus/categorial grammars here}

\paragraph{Kleene Algebra}
Since the early works in the theory of formal languages, Kleene
algebras have played an important role in its development. They
generalize the operations known from regular languages by introducing
operations generalizing language composition, language union and the
Kleene star.  More generally, they are defined as inequational theory
where the inequality is meant to capture language containment. This
theory is extremely successful, having found applications in algebraic
path problems, theory of programming languages, compiler optimizations
and more.

A frequently fruitful research direction is exploring varying
extensions of Kleene algebras, Kleene algebra with tests (KAT) being
one of the most notable ones. Our approach is radically different from
most extensions, which usually aim at modifying or adding new
operations to Kleene algebras, but still keeping it as an inequational
theory. By adopting a category-theoretic treatment and allowing the
``order structure'' to encode more information than merely
inequalities, we were able to extend Kleene algebra to reason about
parsing as well.

\steven{Can we sketch here what an apt categorification of KATs might be? As a
  poor first pass: perhaps a Kleene category with a full involutive
  subcategory that corresponds to the Boolean algebra?.}
\steven{What can we say about internalizing Kleene algebra proofs or extensions?}
\steven{any KA citations needed above?}

%%    - verified imperative implementation: Lambek logic ala separation
%%      logic?
\paragraph{Relation to Separation Logic}

\theoryabbv is similar in spirit to separation logic. Semantically,
they are closely related: linear types in \theoryabbv denote families
of sets indexed by a monoid of strings, whereas separation logic
formulae typically denote families of predicates indexed by an ordered
partial commutative monoid of worlds. The monoidal structure in both
cases is are instances of the category-theoretic notion of Day
convolution monoidal structure. From a separation-logic perspective,
our notion of memory is very primitive: a memory shape is just a
string of characters and the state of the memory is never allowed to
evolve.

Syntactically, our type theory can be viewed as a kind of ``separation
type theory'' in that we care not just about provability of formulae
but the constructive content of our linear terms, as these act as
parsers. The design of \theoryabbv is based on the $\textrm{LNL}_D$
calculus which explicitly was designed to be a dependent
type-theoretic version of separation logic. Many concepts of
separation logic then have direct analogues in our system. The
characters $c$ are analogous to the ``points-to'' formulae $l \mapsto
x$ in separation logic. Our $\otimes,\lto,\tol$ are analogous to the
separating conjunction and magic-wand. The analogue of the persistence
modality $\square$ of separation logic is $\bigoplus_{\_:\uparrow A} I$,
which ``zeros out'' all but the parses from the empty string. If we
think of the portions of the string we are parsing as our notion of
resource, then these make sense as the ``resource-free'' elements of
the type $A$.

This semantic connection to separation logic also suggests an avenue
of future work: rather than working with a dependent type theory, we
could instead make a program logic based on non-commutative separation
logic for verifying imperative implementations of parsers. This could
be accomplished by modifying an existing separation logic
implementation or embedding the logic within \theoryabbv.

%% \paragraph{Linearity with Dependent Types}
%% \steven{I don't know if it worth discussing the different influences of
%%   categorical model here? We take Neel's model and add to it as necessary}

% Vakar might be worth citing to discuss different characterizations of the
% notion of model \cite{vakarSyntaxSemanticsLinear2015}

% \cite{fu2023twolevellineardependenttype} give a compiler for a two-level
% linear dependent type theory. This could serve as a template for a
% compiler implementation of \theoryname that is not just inference rules
% embedded in Agda.

\paragraph{Implementations in Parsing}
\steven{Talk about future implementation-side work}

\cite{fu2023twolevellineardependenttype} give a compiler for a two-level
linear dependent type theory. This could serve as a template for a
compiler implementation of \theoryname that is not just inference rules
embedded in Agda.

\paragraph{Type Checking and Elaboration}
Our focus in this work has been on the verification of parsers for
grammars over strings, but because \theoryabbv allows for the
definition of arbitrarily powerful grammars, the system could also be
used for more sophisticated semantic analysis such as scope checking
or type checking. Alternatively, we could more directly encode type
type systems as linear types in a modified version of \theoryabbv
whose semantics is not in families of sets indexed by \emph{strings}
but families of sets indexed by \emph{abstract syntax trees}.

%% While parsing typically refers to the generation of semantic
%% objects from string input, many tasks in programming can be
%% viewed as parsing of objects with more structure, such as
%% trees with binding structure or graphs. Fundamental to the
%% frontend of many
%% programming language implementations are type systems. In
%% particular, \emph{type checking}
%% --- analogous to language recognition --- and \emph{typed
%%   elaboration} --- analogous to parsing --- arise when
%% producing a semantic object subject to some analysis. Just
%% as our string grammars were given as functors from $\String$
%% to $\Set$, we envision adapting the same philosophy
%% to functors from a \emph{category of trees} to $\Set$ to craft a syntax
%% that natively captures typed elaboration. This suggests an
%% unusual sort of bunched type theory, where context extension
%% no longer resembles concatenation of strings but instead
%% takes on the form of tree constructors.

% type system              ~~ Formal Grammar
% typing derivation        ~~ parse tree
% algorithmic type system  ~~ LL(1), LR(1) grammar
% Uniqueness of derivations ~~ unambiguous grammar
% type elaborator           ~~ semantic actions

\bibliographystyle{plain}
\bibliography{refs.bib}

\appendix

\section{Elided Syntax}

\max{TODO: add prose}

\begin{figure}
  \begin{mathpar}
    \inferrule{~}{\ctxwff \cdot}
    \and
    \inferrule{\ctxwff \Gamma \\ \ctxwffjdg \Gamma X}{\ctxwff {\Gamma, x : X}}
    \inferrule{\Gamma \vdash X : U_i}{\ctxwffjdg \Gamma X}

    \boxed{\linctxwffjdg \Gamma A}

    \inferrule{~}{\linctxwff \Gamma \cdot}
    \and
    \inferrule{\linctxwff \Gamma \Delta \\ \linctxwffjdg \gamma A}{\linctxwff \Gamma {\Delta, a : A}}
  \end{mathpar}
  \caption{Context Rules}
\end{figure}


\max{TODO: add rules for non-linear types $\Pi,\Sigma$, equality etc}

\begin{figure}
  \begin{mathpar}
  \boxed{\Gamma \vdash M : X}

  \inferrule{~}{\Gamma, x : X, \Gamma' \vdash x : X}
  %
  \and
  %
  \inferrule{\Gamma \vdash M : Y \quad \ctxwffjdg \Gamma {X \equiv Y}}{\Gamma \vdash M : X}
  %
  \and
  %
  \inferrule{~}{\Gamma \vdash () : 1}
  %
  \and
  %
  \inferrule{\Gamma \vdash M : X \\ \Gamma \vdash N : \subst Y M x}{\Gamma \vdash (M, N) : \SigTy x X Y}
  %
  \and
%
  \inferrule{\Gamma \vdash M : \SigTy x X Y}{\Gamma \vdash \pi_1\, M : X}
  %
  \and
  %
  \inferrule{\Gamma \vdash M : \SigTy x X Y}{\Gamma \vdash \pi_2\, M : \subst Y {\pi_1\, M} x}
  \and
  \inferrule{\Gamma, x : X \vdash M : Y}{\Gamma \vdash \lamb x M : \PiTy x X Y}
  %
  \and
  %
  \inferrule{\Gamma \vdash M : \PiTy x X Y \\ \Gamma \vdash N : X}{\Gamma \vdash \app M {N} : \subst Y {N} x}
  %
  \and
  %
  \inferrule{\Gamma \vdash M \equiv N : X}{\Gamma \vdash \mathsf{refl} : M =_X N}
  \and
  \textrm{TODO: inductive types (boolean, nat) with their eliminators allowed in }
  \end{mathpar}
  \caption{Intuitionistic typing}
  \label{fig:inttyping}
\end{figure}

\begin{figure}
  \begin{mathpar}
    %
    \inferrule{\Gamma ; a : A , \Delta \vdash e : B}{\Gamma ; \Delta \vdash \lambtol a e : B\tol A}
    \and
    \inferrule{\Gamma ; \Delta \vdash e : B \tol A \\ \Gamma ; \Delta' \vdash e' : A}{\Gamma ; \Delta', \Delta \vdash \apptol e {e'} : B}
    %
    \\

  \end{mathpar}
  \caption{Linear Terms, extended}
\end{figure}

\begin{figure}
  \begin{mathpar}
    \boxed{\Gamma \vdash M \equiv N : X}

    \inferrule{\Gamma \vdash p : M =_X N}{\Gamma \vdash M \equiv N : X}
%
    \and
%
    \inferrule{~}{\Gamma \vdash \app {(\lamb x M)} {N} \equiv \subst x M {N} : X}
%
    \and
%
    \inferrule{~}{\Gamma \vdash m \equiv \lamb x {\app M x} : \PiTy x X Y}
%
    \and
%
    \inferrule{~}{\Gamma \vdash \pi_1\, (M, N) \equiv M : X}
%
    \and
%
    \inferrule{~}{\Gamma \vdash \pi_2\, (M , N) \equiv N : \subst x {M} Y}
%
    \and
%
    \inferrule{~}{\Gamma \vdash M \equiv (\pi_1\, M, \pi_2\, M) : \SigTy x X Y}
%
    \and
%
    % \inferrule{~}{}
%
    \inferrule{~}{\Gamma \vdash M \equiv N : 1}
%
    % \and
%
    % \inferrule{~}{\Gamma \vdash G\, (G^{-1} \, t) \equiv t : G A}

    \\

    \boxed{\Gamma ; \Delta \vdash e \equiv e' : A}

    % \inferrule{~}{\Gamma; \cdot \vdash G^{-1}\, (G \, t ) \equiv t: A}
%
%
    \inferrule{~}{\Gamma; \Delta \vdash \app {(\lamblto a e)} {e'} \equiv \subst e {e'} a : C}
%
    \and
%
    \inferrule{~}{\Gamma; \Delta \vdash e \equiv \lamblto a {\app e a} : A \lto B}
%
    \and
%
    \inferrule{~}{\Gamma; \Delta \vdash \app {(\lambtol a e)} {e'} \equiv \subst e {e'} a : C}
%
    \and
%
    \inferrule{~}{\Gamma; \Delta \vdash e \equiv \lambtol a {\app e a} : B \tol A}
%
    \and
%
    \inferrule{~}{\Gamma; \Delta \vdash \app {(\dlamb x e)} {M} \equiv \subst {e} M x : C}
%
    \and
%
    \inferrule{~}{\Gamma; \Delta \vdash e \equiv \dlamb x {\app e x} : \LinPiTy x X A}
%
    \and
%
%     \inferrule{~}{\Gamma; \Delta \vdash e \equiv e' : \top}
% %
%     \and
%
%     \inferrule{~}{\Gamma; \Delta \vdash e_i \equiv \pi_i (e_1, e_2) : A_i}
% %
%     \and
% %
%     \inferrule{~}{\Gamma; \Delta \vdash e \equiv (\pi_1 e, \pi_2 e) : A\& B}
%
    \and
%
    \inferrule{~}{\Gamma; \Delta \vdash \letin {()} {()} e \equiv e : C}
    %
    \and
    %
    \inferrule{~}{\Gamma; \Delta \vdash \letin {()} e {\subst {e'} {()} a} \equiv \subst {e'} e a : C}
%
    \and
%
    \inferrule{~}{\Gamma; \Delta \vdash \letin {a \otimes b} {e \otimes e'} e'' \equiv e'' \{ e/a , e'/b \} : C}
%
    \and
    %
    \inferrule{~}{\Gamma; \Delta \vdash \letin {a \otimes b} e {\subst {e'} {a \otimes b} c} \equiv \subst {e'} e c : C}
%
    \and
%
    \inferrule{~}{\Gamma;\Delta \vdash \letin {(x, a)} {(M, e)} {e'} \equiv e' \{ M/x , e/a \} : C}
    %
    \and
    %
    \inferrule{~}{\Gamma; \Delta \vdash \letin {(x, a)} e {\subst {e'} {(x, a)} y} \equiv \subst {e'} e y : C}
    \and
    %
    \inferrule
    {~}
    {\Gamma ; \Delta \vdash \equalizerpi {\equalizerin {e}} \equiv e : A}
    %
    \and
    %
    \inferrule
    {~}
    {\Gamma ; \Delta \vdash \equalizerin {\equalizerpi {e}} \equiv e : A}
\end{mathpar}
  \caption{Judgmental equality for linear terms}
  \label{fig:jdgeq}
\end{figure}
\steven{Do we need inverse for $\ltonl A$ in judgmental equality figure?}
\steven{Do we need to say anything about the equalizer proofs being unique in
  judgmental equality figure?}


\end{document}
